% !TeX program = xelatex
\iffalse
\documentclass[11pt]{article}
\usepackage[a4paper,margin=1in]{geometry}
\usepackage{amsmath,amssymb,amsthm,mathtools}
\usepackage{mathrsfs}
\usepackage{booktabs,longtable,tabularx,array}
\usepackage{enumitem}
\usepackage{xcolor}
\usepackage[authoryear,round]{natbib}
\usepackage[hidelinks]{hyperref}
\usepackage{microtype}

\newtheorem{definition}{Definition}[section]
\newtheorem{assumption}[definition]{Assumption}
\newtheorem{theorem}[definition]{Theorem}
\newtheorem{proposition}[definition]{Proposition}
\newtheorem{corollary}[definition]{Corollary}
\newtheorem{lemma}[definition]{Lemma}
\theoremstyle{remark}
\newtheorem{remark}[definition]{Remark}
\newtheorem{example}[definition]{Example}

\newcommand{\E}{\mathbb E}
\newcommand{\Pp}{\mathbb P}
\newcommand{\Qq}{\mathbb Q}
\newcommand{\R}{\mathbb R}
\newcommand{\dd}{\,\mathrm d}
\newcommand{\sgn}{\operatorname{sgn}}
\newcommand{\supp}{\operatorname{supp}}
\newcommand{\Law}{\operatorname{Law}}
\newcommand{\Var}{\operatorname{Var}}
\newcommand{\one}{\mathbf 1}
\newcommand{\classD}{\mathrm{class}\text{-}(D)}
\newcommand{\Normal}{\mathcal N}
\newcommand{\Gam}{\operatorname{Gamma}}
\newcommand{\id}{\mathrm I}
\newcommand{\Leb}{\operatorname{Leb}}
\newcommand{\target}{m}
\renewcommand{\arraystretch}{1.16}
\setlength{\tabcolsep}{4pt}
\setlist{nosep,leftmargin=2em}

\title{Copula-Based Mean-Reverting Markov Processes:\\A Corrected and Self-Contained Theoretical Framework}
\author{}
\date{Revised 15 August 2026}

\begin{document}
\maketitle

\begin{abstract}
This document replaces the theoretical part of
\texttt{meanreversion\_article\_main.tex}.  It separates exact finite-horizon
conditional-mean identities from infinitesimal generator statements, formulates
the absolutely continuous Markov copulas needed in the empirical analysis
directly through their densities, and states the domain,
support, boundary, integrability, and monotonicity conditions that are needed to
turn formal PDE solutions into valid Markov models with observable transition
densities.  The corrected framework retains the valid OU, CIR, monotone-transform,
mixed-copula reset, and nonlocal bond-pricing results.  It treats the two
Gaussian--Gamma calculations as necessary-condition diagnostics, restricts the
quadratic CIR transform to a finite horizon, identifies the admissibility problem
for Kummer branches, excludes the even OU transform from the one-dimensional
monotone framework, and replaces the erroneous second CIR change of variables.
The final section states exactly how each empirical specification relates to the
corrected theory.
\end{abstract}

\section{Probability space, supports, and notation}

Work on a filtered probability space
$(\Omega,\mathcal F,(\mathcal F_t)_{t\ge0},\Pp)$ satisfying the usual
conditions.  Conditional laws given a point are always versions of regular
conditional probability kernels in the sense of
\citet[Chapter~8]{Kallenberg2021}.  Such kernels are probabilistically unique
only up to null sets in the conditioning variable.  Equalities involving an
unspecified version are therefore asserted for the relevant marginal law almost
everywhere; an explicitly selected pointwise version, or continuity of both
sides, is needed before an equality is asserted at every interior state.

For a distribution function $F$, its generalized quantile is
\begin{equation}
Q(u)=F^{-1}(u):=\inf\{x\in\R:F(x)\ge u\},\qquad 0<u<1.
\label{rev:quantile}
\end{equation}
Quantiles are not assumed finite or differentiable at $u=0,1$.  The support of
the law at time $t$ is denoted by $E_t=\supp(F_t)$.  A statement that a CDF is
strictly increasing means strictly increasing on the interior of its support,
not on all of $\R$.  The symbol $m$ denotes the target mean, $\bar r$ a base
diffusion's long-run level, and $s$ a Gamma scale.  These objects were all denoted
by $\theta$ in the original manuscript and must not be conflated.

\section{Markov--copula foundations}

\subsection{Absolutely continuous Markov copulas}

Let $I=(0,1)$ and let $\Leb$ be Lebesgue probability measure on $I$.
Set $\Delta=\{(s,t):0\le s<t\}$.
We use the standard copula and star-product terminology of
\citet{Nelsen2006} and \citet{DarsowNguyenOlsen1992}, subject to the
pointwise-kernel qualifications below.

\begin{assumption}[Absolutely continuous Markov--copula density evolution]
\label{rev:ass-copula-density}
There is a \emph{selected transition-density version}, namely a jointly
Borel-measurable function
$c:\Delta\times I^2\to[0,\infty)$; write
$c_{s,t}(u,v)=c(s,t,u,v)$.  For every $s<t$,
\begin{align}
\int_I c_{s,t}(u,v)\dd v&=1,
&&\text{for every }u\in I,\notag\\
\int_I c_{s,t}(u,v)\dd u&=1,
&&\text{for $\Leb$-a.e. }v\in I,
\label{rev:copula-density-margins}\\
c_{s,t}(u,v)&=\int_Ic_{s,r}(u,w)c_{r,t}(w,v)\dd w,
&&\begin{gathered}
0\le s<r<t,\quad u\in I,\\[-2pt]
\text{for $\Leb$-a.e. }v\in I.
\end{gathered}
\label{rev:copula-density-CK}
\end{align}
For each fixed $(s,r,t,u)$ in \eqref{rev:copula-density-CK}, the exceptional
Lebesgue-null set in $v$ may depend on $(s,r,t,u)$.  For $B\in\mathcal B(I)$,
define the selected kernel
\begin{equation}
K_{s,t}(u,B):=\int_Bc_{s,t}(u,v)\dd v,
\qquad u\in I,\quad s<t,
\label{rev:selected-copula-kernel}
\end{equation}
and put $K_{t,t}(u,B)=\one_B(u)$.  Thus ``for every $u$'' below refers to this
chosen pointwise kernel version; it is stronger than merely specifying a
two-dimensional density up to $\Leb^2$-null sets.
For $s<t$, define
\begin{equation}
C_{s,t}(x,y)=\int_0^x\int_0^y c_{s,t}(u,v)\dd v\dd u,
\qquad (x,y)\in[0,1]^2.
\label{rev:copula-from-density}
\end{equation}
On the diagonal set $C_{t,t}(x,y)=\min(x,y)$ and use the identity transition
operator; no Lebesgue density $c_{t,t}$ is asserted.
\end{assumption}

\begin{proposition}[Density consistency and the copula star product]
\label{rev:prop-density-copula-consistency}
Under Assumption~\ref{rev:ass-copula-density},
\eqref{rev:copula-from-density} is a copula with density $c_{s,t}$.
The kernels \eqref{rev:selected-copula-kernel} satisfy, for every $u\in I$ and
every $B\in\mathcal B(I)$,
\begin{equation}
K_{s,t}(u,B)=\int_IK_{r,t}(w,B)K_{s,r}(u,\dd w),
\qquad 0\le s<r<t.
\label{rev:copula-kernel-CK}
\end{equation}
Consequently, the operators
\begin{equation}
(P_{s,t}g)(u)=\int_Ig(v)c_{s,t}(u,v)\dd v,
\qquad P_{t,t}=\mathrm I,
\label{rev:copula-density-operator}
\end{equation}
form a pointwise Markov evolution preserving the uniform law.  In particular,
the pair-copulas obey the star-product consistency relation
\begin{equation}
C_{s,t}=C_{s,r}*C_{r,t}.
\label{rev:star-product}
\end{equation}
More explicitly, if $C,D$ have selected densities $c_C,c_D$, respectively,
the density-based definition used here is
\begin{equation}
(C*D)(x,y)
=\int_I
 \left\{\int_0^x c_C(a,w)\dd a\right\}
 \left\{\int_0^y c_D(w,b)\dd b\right\}\dd w.
\label{rev:star-product-density-version}
\end{equation}
This equals the familiar formula
$\int_I\partial_2C(x,w)\partial_1D(w,y)\dd w$ when the ordinary partial
derivatives are interpreted as measurable almost-everywhere representatives in
the integration variable $w$.
For every $0\le t_0<\cdots<t_n$, the finite-dimensional density
\begin{equation}
f_{t_0,\ldots,t_n}(u_0,\ldots,u_n)
=\prod_{j=1}^{n}c_{t_{j-1},t_j}(u_{j-1},u_j)
\label{rev:copula-fdd-density}
\end{equation}
is projectively consistent.  Hence these densities define a Markov process
$U$ with uniform marginals on the continuous time index.  No sample-path
continuity is implied by this construction alone.
\end{proposition}

\begin{proof}
Nonnegativity and the two identities in
\eqref{rev:copula-density-margins} give groundedness, two-increasingness, and
uniform margins.  For fixed $(s,r,t,u)$, integrate
\eqref{rev:copula-density-CK} over $B$ and use Tonelli's theorem to obtain
\eqref{rev:copula-kernel-CK}; hence $P_{s,r}P_{r,t}=P_{s,t}$ pointwise.
Another application of Tonelli gives
\eqref{rev:star-product-density-version} and \eqref{rev:star-product}.
The first marginalization in \eqref{rev:copula-fdd-density} uses the column
normalization, the last uses the row normalization, and removal of an
intermediate time uses \eqref{rev:copula-density-CK}.  Thus the family is
projectively consistent.  Since $I$ is a standard Borel space, the
Kolmogorov extension theorem supplies the stated process
\citep[Chapter~8]{Kallenberg2021}.  The full integral details concerning
derivatives and null sets are given in
Appendix~\ref{rev:app-copula-kernel-proof}.
\end{proof}

\begin{remark}[Ordinary derivatives and modified partial Dini derivatives]
\label{rev:rem-modified-Dini}
For the selected density version, define the pointwise conditional CDF
\begin{equation}
H_{s,t}^{c}(u,y):=K_{s,t}(u,(0,y])
=\int_0^yc_{s,t}(u,v)\dd v,
\qquad (u,y)\in I\times[0,1].
\label{rev:selected-conditional-CDF}
\end{equation}
For each fixed $y$, the fundamental theorem for Lebesgue integrals gives
\begin{equation}
H_{s,t}^{c}(u,y)=\partial_1C_{s,t}(u,y)
\quad\text{for $\Leb$-a.e. }u.
\label{rev:ordinary-partial-ae}
\end{equation}
The modified partial Dini derivative $\mathfrak D_1C$ of
\citet[Definition~2.1 and Theorems~2.1--2.2]{FangJiangLiuYang2020} selects a
CDF-valued kernel at every source state, including points where the ordinary
partial derivative is absent.  It and $H^c$ are versions of the same regular
conditional law and therefore agree for $\Leb$-a.e. source state, but need not
agree on a source-null set.  If $u\mapsto H^c(u,y)$ is continuous at a given
$u$---for example under an appropriate continuous, dominated density
version---then the difference quotient for $C(\cdot,y)$ converges there and
all three quantities agree at that $u$.  Thus existence of a Radon--Nikodym
density alone justifies \eqref{rev:ordinary-partial-ae}, not an unrestricted
pointwise equality.  This is precisely why the present density-only framework
selects \eqref{rev:selected-copula-kernel} explicitly rather than silently
identifying a density version with an everywhere-defined derivative; see also
the differentiation results in \citet[Chapter~3]{Folland1999}.
\end{remark}

\begin{remark}[Why some statements say ``every $u$'' and others ``a.e. $v$'']
\label{rev:rem-quantifier-convention}
The source coordinate $u$ indexes the selected probability kernel
$K_{s,t}(u,\cdot)$, so its row normalization and kernel
Chapman--Kolmogorov equation are required for every $u$.  The target coordinate
$v$ is a Radon--Nikodym density variable, and densities and their identities are
intrinsically determined only $\Leb$-a.e. in $v$.  Integrating a density
identity that holds for a.e. $v$ gives a measure identity for every Borel set
$B$.  Conversely, equality of two absolutely continuous kernels for every
$B$ identifies their densities only a.e. in $v$.  These are roles of the two
coordinates, not properties of the letters used to name them.

Star-product consistency of copula CDFs by itself recovers the conditional
composition only for a.e. source state; it does not choose a pointwise
transition function.  The distinction, and sufficient regularity conditions
that upgrade it, are made explicit in Theorems~3.1--3.2 of
\citet{FangJiangLiuYang2020}.  Assumption~\ref{rev:ass-copula-density} imposes
the stronger selected-version property directly.
\end{remark}

\begin{remark}[Scope of the density restriction]
This article deliberately restricts attention to absolutely continuous
pair-copulas because every copula used in the empirical likelihood has a
density.  This is a simplifying scope assumption, not a claim that all Markov
copulas are absolutely continuous.  Consistency
\eqref{rev:copula-density-CK}, rather than the validity of isolated bivariate
copulas at selected horizons, is what supports a continuous-time Markov model.
\end{remark}

\subsection{Uniformization and exact conditional-expectation identities}

\begin{assumption}[Regular marginal representation]
\label{rev:ass-marginal}
For every $t$, $F_t$ is continuous and strictly increasing on $E_t^\circ$,
$Q_t$ is finite on $I$, and $\int_I|Q_t(u)|\dd u<\infty$.  The process $U$ has
the transition densities in Assumption~\ref{rev:ass-copula-density}, is
adapted to $(\mathcal F_t)$, and is Markov with respect to this filtration.
Moreover,
\begin{equation}
X_t=Q_t(U_t)\quad\Pp\text{-a.s.}
\label{rev:uniformization}
\end{equation}
\end{assumption}

\begin{theorem}[Exact quantile--density criterion]
\label{rev:thm-exact-density}
Under Assumption~\ref{rev:ass-marginal}, for $s<t$,
\begin{equation}
\E[X_t\mid\mathcal F_s]
=\int_IQ_t(v)c_{s,t}(U_s,v)\dd v
\quad\Pp\text{-a.s.}
\label{rev:conditional-quantile}
\end{equation}
Consequently, the integrable process $X$ is a martingale with respect to
$(\mathcal F_t)$ if and only if
\begin{equation}
Q_s(u)=\int_IQ_t(v)c_{s,t}(u,v)\dd v
\quad\text{for each fixed }s<t\text{ and for }\Leb\text{-a.e. }u.
\label{rev:exact-martingale}
\end{equation}
The exceptional set in $u$ may depend on the time pair $(s,t)$.  If both sides
of \eqref{rev:exact-martingale} have continuous versions in $u$, the equality
extends from a.e. $u$ to every $u\in I$.
\end{theorem}

\begin{proof}
The Markov property of $U$ with respect to $(\mathcal F_t)$ and
\eqref{rev:uniformization} give \eqref{rev:conditional-quantile}.  Since
$Q_s(U_s)=X_s$ and $U_s$ is uniform, the martingale identity
$\E[X_t\mid\mathcal F_s]=X_s$ is equivalent to
\eqref{rev:exact-martingale} up to Lebesgue-null sets.
\end{proof}

\begin{proposition}[Observable transition density]
\label{rev:prop-observable-density}
If, in addition, for every $t$, $F_t$ is absolutely continuous with density
$p_t>0$ on $E_t^\circ$, define, for every $x\in E_s^\circ$,
\begin{equation}
K^X_{s,t}(x,B)
:=\int_I\one_B\!\left(Q_t(v)\right)c_{s,t}\!\left(F_s(x),v\right)\dd v.
\label{rev:observable-transition-kernel}
\end{equation}
This selected kernel is a regular conditional-law version for $X$, and a
corresponding transition-density version is
\begin{equation}
p^X_{s,t}(x,y)
=c_{s,t}\!\left(F_s(x),F_t(y)\right)p_t(y),
\qquad x\in E_s^\circ,\quad y\in E_t^\circ,
\label{rev:observable-transition-density}
\end{equation}
for every selected source state $x$ and Lebesgue-a.e. target state $y$.
As a conditional law, \eqref{rev:observable-transition-kernel} is unique only
for $p_s(x)\dd x$-a.e. $x$.  The exact quantifiers for its density identities
are
\begin{align}
\int_{E_t^\circ}p^X_{s,t}(x,y)\dd y&=1,
&&\text{every }x\in E_s^\circ,\notag\\
\int_{E_s^\circ}p_s(x)p^X_{s,t}(x,y)\dd x&=p_t(y),
&&\text{Lebesgue-a.e. }y,\notag\\
\int_{E_r^\circ}p^X_{s,r}(x,z)p^X_{r,t}(z,y)\dd z
&=p^X_{s,t}(x,y),
&&\begin{gathered}
s<r<t,\quad\text{every }x\in E_s^\circ,\\[-2pt]
\text{Lebesgue-a.e. }y\in E_t^\circ.
\end{gathered}
\label{rev:observable-density-identities}
\end{align}
Equivalently, the Chapman--Kolmogorov identity holds for
$K^X$ at every selected $x$ and every Borel target set.
\end{proposition}

\begin{proof}
The selected conditional law of $U_t$ given $U_s=F_s(x)$ is
$K_{s,t}(F_s(x),\cdot)$, so its pushforward by $Q_t$ is
\eqref{rev:observable-transition-kernel}.  The change of variables
$v=F_t(y)$ gives \eqref{rev:observable-transition-density}.  Row normalization
gives the first identity in \eqref{rev:observable-density-identities}; column
normalization gives marginal propagation; and the substitution $w=F_r(z)$ in
\eqref{rev:copula-density-CK} gives its density Chapman--Kolmogorov identity.
Integrating the latter in $y$ gives the kernel identity for every Borel set.
\end{proof}

\subsection{Generators and the true-martingale qualification}

For bounded continuous $g$, use the transition operators
\eqref{rev:copula-density-operator}.  The generator, martingale-problem, and
localization conventions follow \citet[Chapter~4]{EthierKurtz1986}; the
semimartingale and class--$(D)$ facts used below follow
\citet[Chapters~I--III]{Protter2005}.

\begin{definition}[Right and extended generators]
\label{rev:def-generator}
For a Feller evolution, the right generator at time $t$ is
\begin{equation}
\mathcal A_tg=\lim_{h\downarrow0}
\frac{P_{t,t+h}g-g}{h}
\label{rev:right-generator}
\end{equation}
whenever the limit exists in the chosen Banach-space norm.  An unbounded
time-dependent quantile $q(t,u)$ is said to lie in the time--space
\emph{extended generator domain} if a localized Dynkin formula holds with
drift $(\partial_t+\mathcal A_t)q$ and the required integrability is finite.
\end{definition}

\begin{theorem}[Generator martingale theorem]
\label{rev:thm-generator-martingale}
Let $q(t,u)=Q_t(u)$ lie in the time--space extended generator domain.  Suppose
that, for every $T<\infty$, the stopped family
\[
\{q(\tau,U_\tau):\tau\text{ is a stopping time with }0\le\tau\le T\}
\]
is uniformly integrable, that is, $q(\cdot,U_\cdot)$ is of class--$(D)$ on
every finite horizon.  Then
$X_t=q(t,U_t)$ is a true martingale if and only if
\begin{equation}
(\partial_t+\mathcal A_t)q(t,u)=0
\quad\text{for }\dd t\otimes\Leb(\dd u)\text{-a.e. }(t,u).
\label{rev:generator-martingale}
\end{equation}
Without the class--$(D)$ or an equivalent uniform-integrability condition, the
zero-drift equation generally establishes only a local martingale.
\end{theorem}

\begin{proof}
The extended Dynkin formula makes
$q(t,U_t)-q(0,U_0)-\int_0^t(\partial_s+\mathcal A_s)q(s,U_s)\dd s$
a local martingale.  Equation \eqref{rev:generator-martingale} therefore gives
a local martingale; class--$(D)$ upgrades it to a true martingale.  Conversely,
the finite-variation part of a special-semimartingale decomposition is unique,
so a true martingale in the stated domain has zero drift almost everywhere.
The measure-theoretic details, including the passage from
$\dd t\otimes\Pp$ to $\dd t\otimes\Leb$, are given in
Appendix~\ref{rev:app-generator-martingale-proof}.
\end{proof}

\subsection{Gaussian and diffusion copula generators}

Let $\Phi$ and $\phi$ be the standard Normal CDF and density.

\begin{proposition}[Brownian and deterministic-time-change generators]
\label{rev:prop-gaussian-generator}
For the uniformized Brownian family and $0<s<t$, the pair-copula is Gaussian
with correlation $\rho_{s,t}=\sqrt{s/t}\in(0,1)$ and density
\begin{equation}
c^{G}_{s,t}(u,v)=
\frac{\phi_{\rho_{s,t}}\!\left(\Phi^{-1}(u),\Phi^{-1}(v)\right)}
     {\phi\!\left(\Phi^{-1}(u)\right)\phi\!\left(\Phi^{-1}(v)\right)},
\label{rev:gaussian-copula-density}
\end{equation}
where $\phi_\rho$ is the standard bivariate Normal density with correlation
$\rho$.  The raw standardized Brownian example is defined on $(0,\infty)$.
For $t>0$, its local generator on $C_c^2(I)$ is
\begin{equation}
\mathcal A_tg(u)=
-\frac{z\phi(z)}{t}g'(u)+\frac{\phi(z)^2}{2t}g''(u),
\qquad z=\Phi^{-1}(u).
\label{rev:brownian-generator}
\end{equation}
If Brownian time is replaced by a differentiable, strictly increasing $h$ with
$h(t)>0$, replace $\rho_{s,t}$ in \eqref{rev:gaussian-copula-density} by
$\sqrt{h(s)/h(t)}\in(0,1)$, and
\begin{equation}
\mathcal A_t^hg(u)=\frac{h'(t)}{h(t)}
\left[-z\phi(z)g'(u)+\frac12\phi(z)^2g''(u)\right].
\label{rev:timechange-generator}
\end{equation}
In particular, $h(t)=e^{2at}$ with $a>0$ gives
\begin{equation}
\mathcal A g(u)=a\phi(z)^2g''(u)-2az\phi(z)g'(u).
\label{rev:stationary-gaussian-generator}
\end{equation}
For both clocks,
$\rho_{s,t}=\rho_{s,r}\rho_{r,t}$ for $s<r<t$; Gaussian conditioning therefore
gives the density Chapman--Kolmogorov identity
\eqref{rev:copula-density-CK}.  A specification whose time index includes
$t=0$ must use a clock with $h(0)>0$ rather than $B_t/\sqrt t$ at the origin.
\end{proposition}

\begin{proof}
For $Z_t=B_t/\sqrt t$, It\^o's formula gives
$\dd Z_t=-Z_t(2t)^{-1}\dd t+t^{-1/2}\dd B_t$.  Applying It\^o again to
$U_t=\Phi(Z_t)$ yields \eqref{rev:brownian-generator}.  The deterministic
time change multiplies the local clock by $h'(t)$ and evaluates the Brownian
coefficient at $h(t)$, which gives \eqref{rev:timechange-generator}.
\end{proof}

\begin{proposition}[Uniformized diffusion generator]
\label{rev:prop-diffusion-copula-generator}
Let
\begin{equation}
\dd R_t=b(t,R_t)\dd t+\sqrt{a(t,R_t)}\dd W_t
\label{rev:general-diffusion}
\end{equation}
take values in an interval $(\ell,r)$.  Suppose that its CDF
$F\in C^{1,2}$ in the interior, $p=\partial_xF>0$, and $a$ is locally $C^1$
in $x$.  Write $F_t^R(x)=F(t,x)$ and $Q_t^R=(F_t^R)^{-1}$.  Assume that, for
every $s<t$, selected kernels
\begin{equation}
K^R_{s,t}(x,B)=\int_Bp^R_{s,t}(x,y)\dd y
\label{rev:diffusion-selected-kernel}
\end{equation}
are jointly Borel in $(s,t,x)$, form a Markov evolution for every
$x\in(\ell,r)$ and Borel $B$, and give the entire regular conditional law,
with no additional boundary atom.  Then a selected version of its pair-copula
density is
\begin{equation}
c_{s,t}(u,v)=
\frac{p^R_{s,t}\!\left(Q_s^R(u),Q_t^R(v)\right)}
     {p\!\left(t,Q_t^R(v)\right)}.
\label{rev:diffusion-copula-density}
\end{equation}
for $(u,v)\in I^2$.  As a Radon--Nikodym derivative it is unique only
$\Leb^2$-a.e.; the displayed formula specifies the version used by the
transition kernel.

Assume additionally that the marginal density and coefficients have enough
interior regularity for the classical forward equation
$p_t=-\partial_x(bp)+\tfrac12\partial_{xx}(ap)$, and that the left-boundary
probability flux satisfies
\begin{equation}
b(t,x)p(t,x)-\tfrac12\partial_x\{a(t,x)p(t,x)\}
\longrightarrow0\quad\text{as }x\downarrow\ell
\label{rev:zero-flux}
\end{equation}
vanishes, then the local generator of $U_t=F_t(R_t)$ is
\begin{align}
\mathcal A_tg(u)
={}&\left[a\,p_x+\tfrac12a_xp\right]_{x=Q_t^R(u)}g'(u)
+\frac12\left[ap^2\right]_{x=Q_t^R(u)}g''(u).
\label{rev:diffusion-copula-generator}
\end{align}
\end{proposition}

\begin{proof}
The conditional change of variables $y=Q_t^R(v)$ gives
\eqref{rev:diffusion-copula-density}.  Its row normalization is the
normalization of the selected kernel \eqref{rev:diffusion-selected-kernel}; its
column normalization holds for a.e. target rank and follows from marginal
propagation
$p(t,y)=\int_\ell^r p(s,x)p^R_{s,t}(x,y)\dd x$.
With $z=Q_r^R(w)$ and $\dd w=p(r,z)\dd z$, the physical transition-density
Chapman--Kolmogorov identity gives \eqref{rev:copula-density-CK} for every
selected source rank and a.e. target rank.  Apply the time-dependent It\^o
formula to $F_t^R(R_t)$ and integrate the forward equation from $\ell$ to $x$.
Condition \eqref{rev:zero-flux} removes the boundary constant.  A final It\^o
application to $g(U_t)$ gives \eqref{rev:diffusion-copula-generator}.  The
It\^o and diffusion-forward-equation results used here are standard; see
\citet[Chapters~5--6]{KaratzasShreve1991}.
\end{proof}

\begin{remark}
The preceding interior result covers OU directly.  CIR is degenerate at zero
and must be treated on $(0,\infty)$ with its nonnegative solution, interior
density, and boundary behavior stated separately; a uniformly elliptic
whole-line assumption cannot be invoked for CIR.
\end{remark}

\section{Correct notions of mean reversion}

\subsection{Local, conditional, and auxiliary notions}

\begin{definition}[Forward local drift]
\label{rev:def-forward-drift}
For an integrable real-valued Markov process with transition kernel
$P_{s,t}(x,\dd y)$, its forward local drift is
\begin{equation}
b_X(t,x)=\lim_{h\downarrow0}
\frac{\int yP_{t,t+h}(x,\dd y)-x}{h},
\label{rev:forward-drift}
\end{equation}
whenever a finite limit exists.  The definition is made for a chosen regular
kernel version and for $F_t$-almost every $x\in E_t$; a continuous version may
be extended to all interior states.
\end{definition}

\begin{definition}[One-dimensional Markov mean reversion]
\label{rev:def-local-mr}
An integrable real-valued Markov process $X$ is \emph{locally mean reverting
toward $m$} if
\begin{enumerate}[label=(\roman*)]
\item $\lim_{t\to\infty}\E[X_t]=m$;
\item the forward drift \eqref{rev:forward-drift} exists and
\begin{equation}
(x-m)b_X(t,x)<0\quad(x\ne m),
\qquad b_X(t,m)=0
\label{rev:local-sign}
\end{equation}
whenever the indicated states belong to the support.
\end{enumerate}
The quantity $b_X$ is a drift, not a ``reversion rate.''
\end{definition}

\begin{definition}[Strong conditional-mean reversion]
\label{rev:def-strong-mr}
An integrable adapted process $X$ has \emph{strong conditional-mean reversion}
toward $m$ if a deterministic function $\mathcal R_{s,t}>0$ exists such that
\begin{align}
\E[X_t-m\mid\mathcal F_s]&=\mathcal R_{s,t}(X_s-m),
      &&0\le s<t,\label{rev:strong-mr}\\
\mathcal R_{s,r}\mathcal R_{r,t}&=\mathcal R_{s,t},\qquad
\lim_{t\to\infty}\mathcal R_{s,t}=0.\label{rev:R-properties}
\end{align}
A principal case is
$\mathcal R_{s,t}=\exp\{\int_s^t\gamma(u)\dd u\}$ with continuous
$\gamma(u)<0$.
\end{definition}

\begin{definition}[Finite-horizon and directional classes]
\label{rev:def-auxiliary-classes}
For model diagnostics we use two weaker labels, neither of which is silently
called global mean reversion.
\begin{enumerate}[label=(\roman*)]
\item A \emph{finite-horizon conditional-reversion model} satisfies
\eqref{rev:strong-mr} and is a one-dimensional Markov model only on a stated
window $0\le t\le T$.
\item A \emph{directional-reversion model} has a drift pointing toward a
threshold $c$, but $c$ has not been proved equal to the limiting mean.
\end{enumerate}
An observed projection of a Markov vector is called a \emph{latent-state
extension}; it is not asserted to be a one-dimensional Markov process.
\end{definition}

\begin{proposition}[Strong reversion implies local reversion]
\label{rev:prop-strong-local}
Suppose Definition~\ref{rev:def-strong-mr} holds with
$\mathcal R_{s,t}=\exp\{\int_s^t\gamma(u)\dd u\}$, where $\gamma$ is continuous and
negative.  Then
\begin{equation}
b_X(t,x)=\gamma(t)(x-m).
\label{rev:linear-drift}
\end{equation}
If $X_s$ is integrable, then $\E[X_t]\to m$, so $X$ satisfies
Definition~\ref{rev:def-local-mr}.
\end{proposition}

\begin{proof}
Set $s=t$ and replace $t$ in \eqref{rev:strong-mr} by $t+h$; differentiation
at $h=0+$ gives \eqref{rev:linear-drift}.  Taking expectations in
\eqref{rev:strong-mr} and using \eqref{rev:R-properties} gives the limiting
mean.  Since $\gamma(t)<0$, the sign condition follows.
\end{proof}

\begin{remark}
The converse is false in general: an instantaneous sign condition and an
unconditional limiting mean do not by themselves determine every finite-horizon
conditional mean.  This is why the local and strong definitions are kept
separate.
\end{remark}

\begin{proposition}[Martingale rescaling]
\label{rev:prop-martingale-rescaling}
Let $X$ be an integrable Markov process satisfying
Definition~\ref{rev:def-strong-mr}.  Then
\begin{equation}
M_t=\frac{X_t-m}{\mathcal R_{0,t}}
\label{rev:martingale-rescaling}
\end{equation}
is a Markov martingale.  Conversely, if $M$ is an integrable Markov martingale,
$\mathcal R_{0,t}>0$ is deterministic and multiplicative, and
$X_t=m+\mathcal R_{0,t}M_t$, then $X$ satisfies
\eqref{rev:strong-mr}.  The two processes have the same pair-copulas.  If
$H_t$ is the marginal CDF of $M_t$, their quantiles satisfy
\begin{equation}
H_t^{-1}(u)=\frac{Q_t(u)-m}{\mathcal R_{0,t}}.
\label{rev:rescaled-quantile}
\end{equation}
When the hypotheses of Theorem~\ref{rev:thm-generator-martingale} hold, this is
equivalent to
\begin{equation}
(\partial_t+\mathcal A_t)Q(t,u)
-\gamma(t)\{Q_t(u)-m\}=0.
\label{rev:strong-mr-generator-PDE}
\end{equation}
\end{proposition}

\begin{proof}
Multiplicativity gives
$\mathcal R_{0,t}=\mathcal R_{0,s}\mathcal R_{s,t}$.  Hence
\[
\E[M_t\mid\mathcal F_s]
=\frac{\mathcal R_{s,t}(X_s-m)}{\mathcal R_{0,t}}
=M_s.
\]
The converse follows by reversing this calculation.  At each time the map
between $X_t$ and $M_t$ is strictly increasing and affine, so it preserves
pair-copulas and gives \eqref{rev:rescaled-quantile}.  Applying the generator
martingale theorem to that quantile and differentiating
$\mathcal R_{0,t}$ gives \eqref{rev:strong-mr-generator-PDE}.
\end{proof}

\subsection{Exact copula and generator characterizations}

\begin{theorem}[Exact copula criterion for strong mean reversion]
\label{rev:thm-exact-mr}
Under Assumption~\ref{rev:ass-marginal}, $X$ satisfies
\eqref{rev:strong-mr} if and only if
\begin{equation}
\int_I\{Q_t(v)-m\}c_{s,t}(u,v)\dd v
=\mathcal R_{s,t}\{Q_s(u)-m\}
\label{rev:exact-mr-copula}
\end{equation}
for each fixed $s<t$ and for $\Leb$-almost every $u$; the exceptional set may
depend on $(s,t)$.  If both sides admit continuous versions in $u$, the equality
extends to every $u\in I$.
\end{theorem}

\begin{proof}
By \eqref{rev:conditional-quantile}, the left side of
\eqref{rev:exact-mr-copula}, evaluated at $u=U_s$, equals
$\E[X_t-m\mid\mathcal F_s]$.  Its right side then equals
$\mathcal R_{s,t}(X_s-m)$.  Since $U_s$ is uniform, equality almost surely is
equivalent to the displayed density identity for $\Leb$-almost every $u$.
\end{proof}

\begin{proposition}[Generator sign criterion]
\label{rev:prop-generator-sign}
Suppose $Q(t,u)$ is in the time--space extended generator domain and the local
drift limit can be interchanged with the quantile representation.  Then
\begin{equation}
b_X(t,Q_t(u))=(\partial_t+\mathcal A_t)Q(t,u).
\label{rev:drift-quantile-generator}
\end{equation}
Consequently, after separately imposing the limiting-mean condition,
Definition~\ref{rev:def-local-mr} is equivalent to
\begin{equation}
\sgn\{(\partial_t+\mathcal A_t)Q(t,u)\}
=\sgn\{m-Q_t(u)\}
\label{rev:generator-sign}
\end{equation}
for almost every interior $u$, with zero drift at $Q_t(u)=m$.
\end{proposition}

\begin{proof}
Insert $X_t=Q_t(U_t)$ into Definition~\ref{rev:def-forward-drift} and use the
definition of the time--space extended generator.  Strict monotonicity converts
the physical-state sign condition into \eqref{rev:generator-sign}.
A complete kernel-integral proof, with separate bounded and weighted conditions
for the moving time difference quotient, is given in
Appendix~\ref{rev:app-drift-identity-proof}.
\end{proof}

\subsection{Stationary Gaussian copula and Gamma-margin diagnostics}

Let $Z$ be the stationary standardized OU process
\begin{equation}
\dd Z_t=-aZ_t\dd t+\sqrt{2a}\dd W_t,\qquad a>0,\qquad Z_t\sim\Normal(0,1).
\label{rev:standardized-ou}
\end{equation}
Its copula is Gaussian with correlation $e^{-a(t-s)}$.

\begin{proposition}[Stationary Gaussian-copula sign criterion]
\label{rev:prop-stationary-gaussian}
Let $F$ have a finite mean $m$, let
$g(z)=F^{-1}(\Phi(z))$, and suppose $g$ is strictly increasing,
$C^2$, and in the extended generator domain of \eqref{rev:standardized-ou}.
Let $z_F$ be the unique number satisfying $g(z_F)=m$.  If
\begin{equation}
\sgn\{g''(z)-zg'(z)\}=\sgn(z_F-z)
\quad\text{for every }z\in\R,
\label{rev:stationary-gaussian-sign}
\end{equation}
then $X_t=g(Z_t)$ is locally mean reverting toward $m$.
\end{proposition}

\begin{proof}
It\^o's formula gives local drift
$a\{g''(z)-zg'(z)\}$.  Since $g$ is strictly increasing,
$\sgn(m-g(z))=\sgn(z_F-z)$.  Stationarity gives
$\E[X_t]=m$ at every time.  Integrability and the extended-domain assumption
justify the drift calculation.
\end{proof}

\begin{example}[Stationary Gamma margin: a necessary root is not sufficient]
\label{rev:ex-gamma-stationary}
Let $G_\nu$ be the quantile of $\Gam(\nu,1)$, with $\nu>0$, and set
$Q(u)=sG_\nu(u)$, $s>0$.  Define
\begin{equation}
u_\nu=F_{\nu,1}(\nu),\qquad z_\nu=\Phi^{-1}(u_\nu).
\label{rev:gamma-mean-quantile}
\end{equation}
At the target mean $m=\nu s$, the zero-drift requirement reduces to
\begin{equation}
\Gamma(\nu)\left(\frac e\nu\right)^\nu\phi(z_\nu)-2z_\nu=0.
\label{rev:gamma-necessary-root}
\end{equation}
Equation \eqref{rev:gamma-necessary-root} is only necessary.  The model is
locally mean reverting only if the full condition
\begin{equation}
\sgn\{\mathcal A Q(u)\}=\sgn\{\nu s-Q(u)\}
\quad\text{for almost every }u\in I
\label{rev:gamma-full-sign}
\end{equation}
also holds, with the required endpoint integrability.  No existence claim is
made from the scalar root alone.
\end{example}

\begin{example}[Time-varying Gamma scale]
\label{rev:ex-gamma-timevarying}
Let $Q_t(u)=s(t)G_\nu(u)$, where $s(t)>0$ is differentiable and
$s(t)\to s_\infty\in(0,\infty)$.  Write $z=\Phi^{-1}(u)$ and
\begin{equation}
\mathscr G_\nu(u)=2\phi(z)zG_\nu'(u)-\phi(z)^2G_\nu''(u).
\label{rev:gamma-G}
\end{equation}
For the Gaussian-copula generator \eqref{rev:stationary-gaussian-generator},
the quantile drift is
\begin{equation}
D_t(u)=s'(t)G_\nu(u)-a s(t)\mathscr G_\nu(u).
\label{rev:gamma-time-drift}
\end{equation}
The correct target is $m=\nu s_\infty$, and the necessary and sufficient local
sign test is
\begin{equation}
\sgn D_t(u)=\sgn\{m-s(t)G_\nu(u)\}
\quad\text{for almost every }u\in I.
\label{rev:gamma-time-full-sign}
\end{equation}
If $u_t$ solves $s(t)G_\nu(u_t)=m$, then
\begin{equation}
\frac{s'(t)}{s(t)}
=a\frac{\mathscr G_\nu(u_t)}{G_\nu(u_t)}
\label{rev:gamma-time-necessary-ode}
\end{equation}
is a necessary zero-point equation, not a sufficient global condition.  Its
limit also inherits the stationary root requirement
\eqref{rev:gamma-necessary-root}.
\end{example}

\section{Mean-reverting monotone transforms of diffusions}

\subsection{Construction theorem and observable density}

Let $R$ be a Markov diffusion on an interval $E$,
\begin{equation}
\dd R_t=b(R_t)\dd t+\sigma(R_t)\dd W_t,
\qquad
\mathcal Lh=b h_x+\tfrac12\sigma^2 h_{xx}.
\label{rev:base-diffusion}
\end{equation}

\begin{theorem}[Strict diffusion-transform theorem]
\label{rev:thm-diffusion-transform}
Let $f\in C^{1,2}([0,\infty)\times E)$ and suppose that
$x\mapsto f(t,x)$ is strictly increasing for every relevant $t$.  Let
$\gamma$ be continuous and let
\begin{equation}
\mathcal R_{s,t}=\exp\!\left\{\int_s^t\gamma(u)\dd u\right\}.
\label{rev:R-gamma}
\end{equation}
Assume that
\begin{equation}
(\partial_t+\mathcal L)f(t,x)=\gamma(t)\{f(t,x)-m\}
\label{rev:transform-PDE}
\end{equation}
and, for every $T<\infty$,
\begin{equation}
\E\int_0^T \mathcal R_{0,t}^{-2}
\sigma(R_t)^2f_x(t,R_t)^2\dd t<\infty.
\label{rev:transform-integrability}
\end{equation}
Then $Y_t=f(t,R_t)$ is Markov and
\begin{equation}
\E[Y_t-m\mid\mathcal F_s]=\mathcal R_{s,t}(Y_s-m).
\label{rev:transform-conditional-mean}
\end{equation}
If $\gamma(t)<0$, $\mathcal R_{s,t}\to0$, and the initial first moment is finite,
then $Y$ has strong and local mean reversion toward $m$.

Conversely, if the strictly monotone $C^{1,2}$ transform satisfies
\eqref{rev:transform-conditional-mean} and its semimartingale decomposition is
integrable, uniqueness of the drift part yields \eqref{rev:transform-PDE}.
\end{theorem}

\begin{proof}
It\^o's formula and \eqref{rev:transform-PDE} give
\[
\dd(Y_t-m)=\gamma(t)(Y_t-m)\dd t
+\sigma(R_t)f_x(t,R_t)\dd W_t.
\]
Multiplying by $\mathcal R_{0,t}^{-1}$ removes the drift.  Condition
\eqref{rev:transform-integrability} makes the resulting stochastic integral a
true square-integrable martingale on every finite horizon, proving
\eqref{rev:transform-conditional-mean}.  Strict invertibility transfers the
Markov property from $R$ to $Y$.  The converse follows from uniqueness of the
finite-variation part.  These It\^o-integral and semimartingale-uniqueness steps
are covered, for example, by \citet{KaratzasShreve1991} and
\citet{Protter2005}.
\end{proof}

\begin{proposition}[Transition density under a monotone transform]
\label{rev:prop-transform-density}
Suppose a selected transition kernel of $R$ has density
$p_{s,t}(x_s,x_t)$ for every source state $x_s$, and the observations $y_j$
lie in $f(j,E)$ with unique inverse states
$x_j=f^{-1}(j,y_j)$ for $j=s,t$.  Then
\begin{equation}
q_{s,t}(y_s,y_t)=p_{s,t}(x_s,x_t)
\left|\frac{\partial f^{-1}(t,y_t)}{\partial y_t}\right|
=\frac{p_{s,t}(x_s,x_t)}{|f_x(t,x_t)|}.
\label{rev:transform-density}
\end{equation}
for every selected source observation $y_s$ and Lebesgue-a.e. target
observation $y_t$.  This is a valid interior density only where the support,
uniqueness, and Jacobian conditions hold.  For the strictly increasing
transforms used here, the probability ranks are unchanged.  Hence the
pair-copula CDFs satisfy the pointwise identity
\begin{equation}
C^Y_{s,t}(u,v)=C^R_{s,t}(u,v),
\qquad (u,v)\in[0,1]^2,
\label{rev:transform-copula-CDF}
\end{equation}
and their Radon--Nikodym densities satisfy
\begin{equation}
c^Y_{s,t}(u,v)=c^R_{s,t}(u,v)
\quad\text{for $\Leb^2$-a.e. }(u,v).
\label{rev:transform-copula-density}
\end{equation}
One may, and in the empirical likelihood does, select the same pointwise
density version on both sides.  The copula invariance in
\eqref{rev:transform-copula-CDF} is the standard invariance under strictly
increasing marginal transformations; see \citet[Chapter~2]{Nelsen2006}.
\end{proposition}

\begin{proof}
Condition on the uniquely recovered state $R_s=x_s$ and apply the
one-dimensional change-of-variables theorem at time $t$.  Moreover,
$F_t^Y(f(t,x))=F_t^R(x)$ for every interior $x$.  Substitution into the joint
CDF proves \eqref{rev:transform-copula-CDF}; uniqueness of
Radon--Nikodym derivatives proves \eqref{rev:transform-copula-density}.
\end{proof}

\subsection{OU base process}

Use the centered OU process
\begin{equation}
\dd R_t=-\kappa R_t\dd t+\sigma\dd W_t,
\qquad \kappa>0,\quad\sigma>0.
\label{rev:OU}
\end{equation}
For $\Delta=t-s$,
\begin{equation}
R_t\mid R_s=x\sim
\Normal\!\left(e^{-\kappa\Delta}x,
\frac{\sigma^2}{2\kappa}(1-e^{-2\kappa\Delta})\right).
\label{rev:OU-transition}
\end{equation}
The noncentered OU is obtained by replacing $R_t$ with $R_t-\bar r$.
It therefore satisfies
\begin{equation}
\E[R_t-\bar r\mid\mathcal F_s]
=e^{-\kappa(t-s)}(R_s-\bar r),
\qquad
\lim_{t\to\infty}\E[R_t\mid\mathcal F_s]=\bar r.
\label{rev:OU-conditional-mean}
\end{equation}
Set $\gamma=-d$ with $d>0$.  The transform PDE is
\begin{equation}
f_t+\tfrac12\sigma^2f_{xx}-\kappa xf_x+d(f-m)=0.
\label{rev:OU-transform-PDE}
\end{equation}

\begin{example}[Affine transformed OU]
\label{rev:ex-affine-OU}
For $A_0>0$ and $B_0\in\R$,
\begin{equation}
f(t,x)=m+A_0e^{(\kappa-d)t}x+B_0e^{-dt}
\label{rev:affine-OU-transform}
\end{equation}
solves \eqref{rev:OU-transform-PDE} and is strictly increasing on $\R$.
If the base OU is stationary, then
\begin{equation}
Y_t\sim\Normal\!\left(m+B_0e^{-dt},
\frac{A_0^2\sigma^2}{2\kappa}e^{2(\kappa-d)t}\right).
\label{rev:affine-OU-margin}
\end{equation}
The variance may grow when $\kappa>d$; this does not contradict the first-moment
identity \eqref{rev:transform-conditional-mean}, but the transformed process is
then not stationary.
\end{example}

\begin{proof}
Substitution leaves the coefficient equations
$A'+(d-\kappa)A=0$ and $B'+dB=0$.  The derivative
$f_x=A_0e^{(\kappa-d)t}$ is positive.  The marginal law follows from the affine
Normal transformation.
\end{proof}

\begin{example}[Exponential transformed OU]
\label{rev:ex-exponential-OU}
Let
\begin{align}
f(t,x)&=m+e^{A(t)x+B(t)}+C(t),\label{rev:exp-OU-transform}\\
A(t)&=A_0e^{\kappa t},\label{rev:exp-OU-A}\\
B(t)&=B_0-\frac{\sigma^2A_0^2}{4\kappa}
             (e^{2\kappa t}-1)-dt,\label{rev:exp-OU-B}\\
C(t)&=C_0e^{-dt}.\label{rev:exp-OU-C}
\end{align}
If $A_0>0$, this is strictly increasing and its support is
$(m+C(t),\infty)$.  It is positive-valued when $m+C(t)\ge0$.  For stationary
$R$, the shifted variable $Y_t-m-C(t)$ is lognormal with log mean $B(t)$ and
log variance $A(t)^2\sigma^2/(2\kappa)$.
\end{example}

\begin{proof}
Substitution in \eqref{rev:OU-transform-PDE} gives
\[
A'=\kappa A,\qquad B'=-\tfrac12\sigma^2A^2-d,
\qquad C'=-dC,
\]
whose solution is \eqref{rev:exp-OU-A}--\eqref{rev:exp-OU-C}.  The derivative
$f_x=Ae^{Ax+B}$ is positive.  Gaussian exponential moments verify
\eqref{rev:transform-integrability} on finite horizons.
\end{proof}

\begin{proposition}[OU Kummer candidate is not an admissible monotone model]
\label{rev:prop-OU-Kummer-failure}
Let $J(p,q,z)$ solve
$zJ_{zz}+(q-z)J_z-pJ=0$.  The separated family
\begin{equation}
f(t,x)=m+e^{-\zeta t}
J\!\left(\frac{\zeta-d}{2\kappa},\frac12,
         \frac{\kappa x^2}{\sigma^2}\right)
\label{rev:OU-Kummer}
\end{equation}
solves \eqref{rev:OU-transform-PDE} wherever it is smooth.  Every nonconstant
member is, however, noninjective on the OU state space $\R$ and therefore does
not satisfy Theorem~\ref{rev:thm-diffusion-transform} or
Proposition~\ref{rev:prop-transform-density}.
\end{proposition}

\begin{proof}
Direct substitution gives the stated Kummer equation.  The spatial dependence
is through $x^2$, so $f(t,x)=f(t,-x)$ for every $x$.  A nonconstant even function
cannot be one-to-one on $\R$.
\end{proof}

\begin{remark}
A density formed by summing two inverse branches would define a different
observation model.  The transformation would no longer preserve the base
copula, and the observed scalar process need not be Markov; this is outside the
present one-dimensional monotone framework.
\end{remark}

\subsection{CIR base process}

Let
\begin{equation}
\dd R_t=\kappa(\bar r-R_t)\dd t+\sigma\sqrt{R_t}\dd W_t,
\qquad \kappa>0,\quad\bar r>0,\quad\sigma>0,
\label{rev:CIR}
\end{equation}
with state space $[0,\infty)$.  The Feller condition
$2\kappa\bar r\ge\sigma^2$ is a convenient sufficient condition for zero to be
unattainable when the initial state is positive.  For $\Delta=t-s$, the exact
transition law is
\begin{align}
R_t\mid R_s=x&\overset d=c_\Delta
\chi^{\prime 2}_{\nu}(\lambda_\Delta),\label{rev:CIR-transition}\\
c_\Delta&=\frac{\sigma^2(1-e^{-\kappa\Delta})}{4\kappa},
\qquad
\nu=\frac{4\kappa\bar r}{\sigma^2},
\qquad
\lambda_\Delta=
\frac{4\kappa e^{-\kappa\Delta}x}
{\sigma^2(1-e^{-\kappa\Delta})}.
\end{align}
The conditional mean has the same linear reversion form as OU:
\begin{equation}
\E[R_t-\bar r\mid\mathcal F_s]
=e^{-\kappa(t-s)}(R_s-\bar r).
\label{rev:CIR-conditional-mean}
\end{equation}
When started from its invariant law,
\begin{equation}
R_t\sim\Gam\!\left(\frac{2\kappa\bar r}{\sigma^2},
                    \frac{\sigma^2}{2\kappa}\right),
\label{rev:CIR-stationary}
\end{equation}
where the second parameter is scale.  With $\gamma=-d$, the corrected
transform PDE is
\begin{equation}
f_t+\tfrac12\sigma^2x f_{xx}
+\kappa(\bar r-x)f_x+d(f-m)=0.
\label{rev:CIR-transform-PDE}
\end{equation}
The boundary statement, noncentral chi-square transition law, conditional mean,
and invariant Gamma law are the standard CIR results; see
\citet{CoxIngersollRoss1985}.

\begin{example}[Affine transformed CIR]
\label{rev:ex-affine-CIR}
For $A_0>0$ and $B_0\in\R$, define
\begin{align}
A(t)&=A_0e^{(\kappa-d)t},\label{rev:affine-CIR-A}\\
B(t)&=m-A_0\bar r e^{(\kappa-d)t}+B_0e^{-dt},
\label{rev:affine-CIR-B}\\
f(t,x)&=A(t)x+B(t).
\label{rev:affine-CIR-transform}
\end{align}
Then $f$ solves \eqref{rev:CIR-transform-PDE}, is strictly increasing on the
CIR state space, and has support $[B(t),\infty)$.  The original manuscript's
formula is the special case $m=\bar r$ after a harmless reparameterization of
$B_0$.
\end{example}

\begin{proof}
The affine coefficient equations are
$A'+(d-\kappa)A=0$ and
$B'+d(B-m)+\kappa\bar r A=0$.  Equations
\eqref{rev:affine-CIR-A}--\eqref{rev:affine-CIR-B} solve them, and $f_x=A>0$.
Polynomial moments of CIR verify the finite-horizon martingale condition.
\end{proof}

\begin{proposition}[Quadratic CIR transform: exact finite-horizon condition]
\label{rev:prop-quadratic-CIR}
Let
\begin{equation}
f(t,x)=A(t)x^2+B(t)x+C(t),\qquad A_0>0,
\label{rev:quadratic-CIR-transform}
\end{equation}
and set
\begin{equation}
\eta=\frac{A_0(\sigma^2+2\kappa\bar r)}{\kappa}.
\label{rev:quadratic-eta}
\end{equation}
The coefficient solution is
\begin{align}
A(t)&=A_0e^{(2\kappa-d)t},\label{rev:quadratic-A}\\
B(t)&=e^{(\kappa-d)t}
      \{B_0-\eta(e^{\kappa t}-1)\},\label{rev:quadratic-B}\\
C(t)&=m+\frac{\bar r\eta}{2}e^{(2\kappa-d)t}
-\bar r(B_0+\eta)e^{(\kappa-d)t}
\notag\\
&\quad+
\left(C_0-m+\bar rB_0+\frac{\bar r\eta}{2}\right)e^{-dt}.
\label{rev:quadratic-C}
\end{align}
It solves \eqref{rev:CIR-transform-PDE}.  On a fixed window $[0,T]$, the map
$x\mapsto f(t,x)$ is increasing on $[0,\infty)$ for every $t\le T$ if and only if
\begin{equation}
B_0\ge\eta(e^{\kappa T}-1).
\label{rev:quadratic-monotonicity}
\end{equation}
Under a strict inequality the interior inverse Jacobian is strictly finite.
No finite $B_0$ makes a nondegenerate quadratic transform globally increasing
for all $t\ge0$ under the standard CIR parameter restrictions.
\end{proposition}

\begin{proof}
Matching the coefficients of $x^2,x,1$ gives
\begin{align*}
A'+(d-2\kappa)A&=0,\\
B'+(d-\kappa)B+(\sigma^2+2\kappa\bar r)A&=0,\\
C'+d(C-m)+\kappa\bar r B&=0,
\end{align*}
which yields \eqref{rev:quadratic-A}--\eqref{rev:quadratic-C}.  Since
$f_x=2A(t)x+B(t)$ and $A(t)>0$, monotonicity on the nonnegative half-line is
equivalent to $B(t)\ge0$.  Formula \eqref{rev:quadratic-B} is nonnegative for
all $t\le T$ exactly under \eqref{rev:quadratic-monotonicity}.  Its bracket
tends to $-\infty$ as $t\to\infty$, proving the global impossibility.
\end{proof}

\begin{remark}[Why $\kappa<0$ is not the same CIR model]
Algebraically, when $\kappa<0$ the quadratic map can be globally increasing if
$A_0>0$ and $\min\{B_0,B_0+\eta\}\ge0$.  This does not rescue the standard CIR
model.  If $\bar r>0$, the boundary drift $\kappa\bar r<0$ points out of the
nonnegative state space.  If $\bar r<0$ is chosen to make the intercept positive,
the drift is explosive and there is no invariant Gamma law or mean reversion.
Such a process would be a separately defined supercritical square-root diffusion,
not a relaxation of \eqref{rev:CIR}.
\end{remark}

\begin{corollary}[Quadratic inverse and density]
\label{rev:cor-quadratic-inverse}
Under the strict finite-window condition and for an interior observation
$y>C(t)$, the unique nonnegative inverse is
\begin{equation}
x=\frac{2\{y-C(t)\}}
{B(t)+\sqrt{B(t)^2+4A(t)\{y-C(t)\}}},
\qquad
\left|\frac{\partial x}{\partial y}\right|
=\frac1{2A(t)x+B(t)}.
\label{rev:quadratic-inverse}
\end{equation}
Together with \eqref{rev:CIR-transition}, this gives the finite-horizon
observation density through \eqref{rev:transform-density}.
\end{corollary}

\begin{example}[Exponential transformed CIR]
\label{rev:ex-exponential-CIR}
Let
\begin{align}
f(t,x)&=m+e^{A(t)x+B(t)}+D(t),\label{rev:exp-CIR-transform}\\
A(t)&=\frac{2\kappa}
{\sigma^2+(2\kappa/A_0-\sigma^2)e^{-\kappa t}},
\label{rev:exp-CIR-A}\\
B(t)&=-dt-\frac{2\kappa\bar r}{\sigma^2}
\log\!\left[(2\kappa/A_0-\sigma^2)+\sigma^2e^{\kappa t}\right]
\notag\\
&\quad+B_0+
\frac{2\kappa\bar r}{\sigma^2}\log(2\kappa/A_0),
\label{rev:exp-CIR-B}\\
D(t)&=D_0e^{-dt}.
\label{rev:exp-CIR-D}
\end{align}
For $A_0>0$, the logarithm arguments and denominator are positive and the map
is strictly increasing.  Its support is
\begin{equation}
[m+D(t)+e^{B(t)},\infty).
\label{rev:exp-CIR-support}
\end{equation}
If the base CIR is stationary, a sufficient and necessary initial exponential
moment restriction is
\begin{equation}
0<A_0<\frac{2\kappa}{\sigma^2}.
\label{rev:exp-CIR-moment}
\end{equation}
Under these conditions the positive local-martingale component has constant
expectation and is therefore a true martingale; this verifies the conditional
mean identity even when the square-integrability criterion
\eqref{rev:transform-integrability} is stronger than needed.
\end{example}

\begin{proof}
The coefficient equations are
\[
A'+\tfrac12\sigma^2A^2-\kappa A=0,
\qquad B'+d+\kappa\bar r A=0,
\qquad D'+dD=0.
\]
Their solution is \eqref{rev:exp-CIR-A}--\eqref{rev:exp-CIR-D}.  Since
$R_t\ge0$, the minimum of the exponential term is $e^{B(t)}$, proving
\eqref{rev:exp-CIR-support}.  The stationary Gamma MGF is finite exactly below
rate $2\kappa/\sigma^2$, which yields \eqref{rev:exp-CIR-moment}.  More
explicitly, put $p_0=\sigma^2A_0/(2\kappa)\in(0,1)$ and
$q=2\kappa\bar r/\sigma^2$.  Then
\[
B(t)=B_0-dt-q\log(1-p_0+p_0e^{\kappa t}),
\qquad
1-\frac{\sigma^2A(t)}{2\kappa}
=\frac{1-p_0}{1-p_0+p_0e^{\kappa t}}.
\]
Using the stationary Gamma MGF shows
\[
\E\exp\{dt+A(t)R_t+B(t)\}=e^{B_0}(1-p_0)^{-q},
\]
which is constant.  The nonnegative local-martingale component is therefore a
true martingale.
\end{proof}

\subsection{CIR special functions and the corrected reduction}

Let $M(p,q,z)={}_1F_1(p;q;z)$ denote Kummer's function and $U(p,q,z)$
Tricomi's function, with limiting definitions used at exceptional integer
parameters.  We use the definitions, derivative identities, and asymptotic
conventions of \citet[Chapter~13]{OlverEtAl2010}.

\begin{proposition}[CIR Kummer separated family]
\label{rev:prop-CIR-Kummer}
Set
\begin{equation}
p=\frac{\zeta-d}{\kappa},\qquad
q=\frac{2\kappa\bar r}{\sigma^2},\qquad
z=\frac{2\kappa x}{\sigma^2}.
\label{rev:CIR-Kummer-parameters}
\end{equation}
If
\begin{equation}
J(p,q,z)=C_1M(p,q,z)+C_2U(p,q,z)
\label{rev:CIR-Kummer-branch}
\end{equation}
is a well-defined branch satisfying the relevant boundary conditions, then
\begin{equation}
f(t,x)=m+e^{-\zeta t}J(p,q,z)
\label{rev:CIR-Kummer-transform}
\end{equation}
solves \eqref{rev:CIR-transform-PDE}.  It defines a model in the present
framework only if the selected branch is strictly increasing on $[0,\infty)$,
has a unique stable inverse on its image, and satisfies the moment condition in
Theorem~\ref{rev:thm-diffusion-transform} for the stated initial law.

For example, $p,q>0$, $C_1>0$, $C_2=0$ gives
\begin{equation}
\partial_zM(p,q,z)=\frac pqM(p+1,q+1,z)>0,
\label{rev:Kummer-M-derivative}
\end{equation}
but this branch is not integrable under the stationary Gamma law because of its
$e^z$ asymptotic term.  It may be integrable at each finite time from a
deterministic initial state.  Thus the separated equation is correct, whereas
an unspecified ``special-function CIR model'' is not a complete statistical
model.
\end{proposition}

\begin{proof}
Substitution gives
$zJ_{zz}+(q-z)J_z-pJ=0$.  The derivative identity and standard positive series
for $M$ prove monotonicity in the stated branch.  The cited Kummer asymptotic is
$M(p,q,z)\sim\Gamma(q)e^zz^{p-q}/\Gamma(p)$; multiplying by the invariant
$z^{q-1}e^{-z}$ density leaves a nonintegrable $z^{p-1}$ tail for $p>0$.
\end{proof}

\begin{proposition}[The second CIR transform in the original manuscript is false]
\label{rev:prop-CIR-false-transform}
The exponential change of variables claimed in the original manuscript does
not transform \eqref{rev:CIR-transform-PDE} into $u_\tau=zu_{zz}$; direct
substitution has a nonzero residual for generic parameters.  A correct direct
reduction is
\begin{equation}
f(t,x)-m=e^{-dt}u(\tau,z),\qquad
z=xe^{\kappa t},\qquad
\tau=\frac{\sigma^2}{2\kappa}(1-e^{\kappa t}),
\label{rev:CIR-correct-variables}
\end{equation}
which yields
\begin{equation}
u_\tau=zu_{zz}+\frac{2\kappa\bar r}{\sigma^2}u_z.
\label{rev:CIR-correct-reduction}
\end{equation}
Therefore the three functions listed after the false reduction are not thereby
solutions of the original CIR transform PDE.
\end{proposition}

\begin{proof}
Differentiate \eqref{rev:CIR-correct-variables} and substitute into
\eqref{rev:CIR-transform-PDE}; the $\kappa z u_z$ terms cancel and the remaining
terms give \eqref{rev:CIR-correct-reduction}.  The coefficient of $u_z$ is
strictly positive under the standard CIR parameters and cannot generally be
discarded.
\end{proof}

\section{Mixed copula families as stationary resets}

\subsection{Transition-density construction and preservation of mean reversion}

Let $P_{s,t}$ be a Markov evolution on a state space $E$ with common invariant
probability law $\pi$; thus the kernel Chapman--Kolmogorov identity holds for
every selected source state $x$ and every measurable target set.  Assume that,
for $s<t$, selected jointly measurable versions are dominated by a common
$\sigma$-finite measure $\mu_0$,
\[
P_{s,t}(x,\dd y)=p_{s,t}(x,y)\mu_0(\dd y),
\qquad \pi(\dd y)=p_\pi(y)\mu_0(\dd y).
\]
Here the first equality is imposed for every $x$ as an equality of measures;
the densities in $y$ remain unique only $\mu_0$-a.e.
Define the independence/reset operator
\begin{equation}
\Pi(x,A)=\pi(A),\qquad (\Pi h)(x)=\int_Eh(z)\pi(\dd z).
\label{rev:reset-kernel}
\end{equation}
Let $\lambda:[0,\infty)\to[0,\infty)$ be right-continuous and locally bounded,
and set
\begin{equation}
a_{s,t}=\exp\!\left\{-\int_s^t\lambda(u)\dd u\right\}.
\label{rev:survival-weight}
\end{equation}

\begin{theorem}[Consistent reset density mixture]
\label{rev:thm-mixed-kernel}
The kernels
\begin{equation}
P^{\mathrm{mix}}_{s,t}
=a_{s,t}P_{s,t}+(1-a_{s,t})\Pi
\label{rev:mixed-kernel}
\end{equation}
form a Markov evolution and preserve $\pi$.  Their transition densities are
\begin{equation}
p^{\mathrm{mix}}_{s,t}(x,y)
=a_{s,t}p_{s,t}(x,y)+(1-a_{s,t})p_\pi(y).
\label{rev:mixed-transition-density}
\end{equation}
where the right-hand side defines the selected version; its density
Chapman--Kolmogorov identity holds for every $x$ and $\mu_0$-a.e. $y$, while
the corresponding kernel identity holds for every $x$ and every measurable
target set.
Assume additionally that $E\subseteq\R$, the process is started from $\pi$,
the invariant CDF $F_\pi$ is continuous and strictly increasing on the
interior of its support, and the stationary base pair-copula has density
$c_{s,t}$.  Then the mixed pair-copula and its density are
\begin{equation}
C^{\mathrm{mix}}_{s,t}
=a_{s,t}C_{s,t}+(1-a_{s,t})\Pi_C,
\qquad \Pi_C(u,v)=uv,
\label{rev:mixed-copula}
\end{equation}
\begin{equation}
c^{\mathrm{mix}}_{s,t}(u,v)
=a_{s,t}c_{s,t}(u,v)+(1-a_{s,t}).
\label{rev:mixed-copula-density}
\end{equation}
The copula identity \eqref{rev:mixed-copula} is pointwise on $[0,1]^2$.
Equation \eqref{rev:mixed-copula-density} is a $\Leb^2$-a.e. density identity,
and its right-hand side is the selected version used by the mixed kernel; with
that selection, its density consistency holds for every source rank $u$ and
a.e. target rank $v$.
If the base generator is $\mathcal G_t$, then, on any invariant generator
domain whose functions are integrable under $\pi$,
\begin{equation}
\mathcal G_t^{\mathrm{mix}}h(x)
=\mathcal G_th(x)+\lambda(t)
\left\{\int_Eh(z)\pi(\dd z)-h(x)\right\}.
\label{rev:mixed-generator}
\end{equation}
\end{theorem}

\begin{proof}
The survival weights are multiplicative.  Invariance gives
$P\Pi=\Pi P=\Pi$ and $\Pi^2=\Pi$.  Expanding
$P^{\mathrm{mix}}_{s,r}P^{\mathrm{mix}}_{r,t}$ therefore reduces it to
\eqref{rev:mixed-kernel} with weight $a_{s,r}a_{r,t}=a_{s,t}$.
Taking densities gives \eqref{rev:mixed-transition-density}; uniformization
under stationary start gives \eqref{rev:mixed-copula} and
the a.e. identity \eqref{rev:mixed-copula-density}.  Selecting its right-hand
side as the density version, its margins equal one, and direct expansion
using \eqref{rev:copula-density-margins}--\eqref{rev:copula-density-CK} gives
\[
\int_Ic^{\mathrm{mix}}_{s,r}(u,w)c^{\mathrm{mix}}_{r,t}(w,v)\dd w
=a_{s,t}c_{s,t}(u,v)+(1-a_{s,t})
=c^{\mathrm{mix}}_{s,t}(u,v).
\]
for every $u$ and a.e. $v$; integration in $v$ gives exact kernel consistency
for every target Borel set.
The right derivative of \eqref{rev:mixed-kernel} at $t=s+$ gives
\eqref{rev:mixed-generator}.
\end{proof}

\begin{theorem}[Exact strengthening of conditional mean reversion]
\label{rev:thm-mixed-mr}
Suppose the base process has invariant mean
$m=\int_Ez\pi(\dd z)$ and satisfies
\begin{equation}
\E[Z_t-m\mid Z_s=x]=\mathcal R_{s,t}(x-m).
\label{rev:base-mr-for-mix}
\end{equation}
Then the reset-mixture process satisfies
\begin{equation}
\E[Y_t-m\mid Y_s=x]=a_{s,t}\mathcal R_{s,t}(x-m).
\label{rev:mixed-exact-mr}
\end{equation}
Its local drift is
\begin{equation}
b_{\mathrm{mix}}(t,x)=b_0(t,x)+\lambda(t)(m-x).
\label{rev:mixed-local-drift}
\end{equation}
Thus a nonnegative reset intensity preserves the local sign condition and, for
$\lambda>0$, strictly accelerates the decay of conditional mean deviations
away from $m$.
\end{theorem}

\begin{proof}
The base component of \eqref{rev:mixed-kernel} contributes
$a_{s,t}\mathcal R_{s,t}(x-m)$.  A draw from $\pi$ has centered mean zero, so the reset
component contributes zero.  This proves \eqref{rev:mixed-exact-mr}; its
right derivative gives \eqref{rev:mixed-local-drift}.
\end{proof}

\begin{remark}[A precise theoretical advantage]
The mixed model does not alter any one-time marginal under stationary start,
but multiplies the base conditional-mean persistence by $a_{s,t}$.  Hence
$a_{s,t}\mathcal R_{s,t}$, the implied half-life, or the additive rate
$-\partial_{t+}\log a_{s,t}=\lambda(t)$ are clear structural measures of its
extra reversion channel.  Whether this produces superior out-of-sample log
scores or calibration is a separate empirical question.
\end{remark}

\subsection{Poisson realization}

\begin{theorem}[Independent-copy reset construction]
\label{rev:thm-poisson-construction}
Let $\{Z^{(j)}:j\ge0\}$ be independent copies of a stationary Markov process
with transition evolution $P$ and invariant law $\pi$.  Let $N$ be an
independent nonhomogeneous Poisson process with intensity $\lambda(t)$, and set
\begin{equation}
Y_t=Z_t^{(N_t)}.
\label{rev:poisson-copy-process}
\end{equation}
Then $Y$ is Markov, stationary in its one-time marginals, and has transition
kernel \eqref{rev:mixed-kernel} and transition density
\eqref{rev:mixed-transition-density}.  Under the stationary continuous-margin
conditions of Theorem~\ref{rev:thm-mixed-kernel}, it also has Copula density
\eqref{rev:mixed-copula-density}.  A constant $\lambda$ gives
$a_{s,t}=e^{-\lambda(t-s)}$; a positive deterministic seasonal intensity is
also admissible.
\end{theorem}

\begin{proof}
Conditional on no jump in $(s,t]$, whose probability is $a_{s,t}$, the same
copy evolves with kernel $P_{s,t}$.  Conditional on at least one jump, the
active copy at $t$ is independent of the history through time $s$ and has law
$\pi$.  The resulting conditional law is \eqref{rev:mixed-kernel}, whose
density is \eqref{rev:mixed-transition-density}, and it is a
Markov evolution by Theorem~\ref{rev:thm-mixed-kernel}.
\end{proof}

\section{Risk-neutral bond pricing with reset risk}

This section is a pricing theorem under a risk-neutral measure $\Qq$, not a
claim that a physical-measure stablecoin spot series identifies bond prices.

\begin{theorem}[Nonlocal Feynman--Kac equation]
\label{rev:thm-bond-PDE}
Under $\Qq$, let the short rate be a stationary-reset OU process with generator
\begin{equation}
\mathcal Gh(y)=\kappa(m-y)h'(y)+\tfrac12\sigma^2h''(y)
+\lambda\left\{\int_\R h(z)\pi(\dd z)-h(y)\right\},
\label{rev:bond-generator}
\end{equation}
where $\pi=\Normal(m,\sigma^2/(2\kappa))$.  Suppose the exponential discount
functional is integrable and that the classical Feynman--Kac hypotheses hold;
alternatively, assume conditions giving a unique viscosity solution with the
stated growth class.  For jump generators and the corresponding nonlocal
Feynman--Kac equations, see \citet[Chapters~8 and~12]{ContTankov2004}.  Then
\begin{equation}
g(t,y)=\E^{\Qq}\!\left[
\exp\!\left\{-\int_t^TY_s\dd s\right\}\middle|Y_t=y\right]
\label{rev:bond-value}
\end{equation}
solves
\begin{equation}
\begin{cases}
g_t+\kappa(m-y)g_y+\tfrac12\sigma^2g_{yy}
+\lambda\{\int_\R g(t,z)\pi(\dd z)-g(t,y)\}-yg=0,\\
g(T,y)=1.
\end{cases}
\label{rev:bond-PDE}
\end{equation}
\end{theorem}

\begin{proof}
Theorem~\ref{rev:thm-mixed-kernel} gives the generator
\eqref{rev:bond-generator}.  Apply the Feynman--Kac formula with killing rate
$y$ and terminal payoff one.  Rearranging the generator equation yields
\eqref{rev:bond-PDE}.
\end{proof}

\section{Embedding of the empirical specifications}

The following model identities are part of the theory--empirics interface.
They prevent four distinct notions---a coded forward map, an invertible map, a
normalized transition density, and a successfully estimated structural
model---from being collapsed into one ``implemented'' flag.

Use the following classes:
\begin{description}
\item[S] one-dimensional structural model covered by the corrected theorems;
\item[C] covered only after an explicit condition, a finite horizon, or a
change in the observed state;
\item[L] legitimate latent-state or horizon-specific predictive extension,
but not a one-dimensional continuous-time Markov model for the observation;
\item[B] benchmark not carrying the structural claim;
\item[N] not embeddable under its current model identity.
\end{description}

\scriptsize
\begin{longtable}{p{0.035\textwidth}p{0.205\textwidth}p{0.045\textwidth}p{0.635\textwidth}}
\toprule
No. & Empirical registry name & Class & Correct theoretical interpretation\\
\midrule
\endfirsthead
\toprule
No. & Empirical registry name & Class & Correct theoretical interpretation\\
\midrule
\endhead
1 & OUF & S & Exact OU transition; $\kappa,\sigma>0$ and finite first moment give strong reversion.\\
2 & ThresholdOU\_\allowbreak{}Heteroskedastic & B & The drift points to a threshold, but with unequal $\kappa_\pm$ the invariant mean is not that threshold.  The current Euler Gaussian QMLE is not the exact threshold-diffusion transition.\\
3 & CIR\_Price & S & Standard CIR in theory.  Any Normal moment-matching numerical fallback must be disclosed as an approximate likelihood rather than an exact CIR density.\\
4 & CIR\_\allowbreak{}DepegPressure & C & Valid only when the absolute depeg amplitude itself is defined as the CIR state.  The sign of the price deviation is lost; a signed-price model requires an augmented sign state.\\
5 & AffineTransformedOU & S & Example~\ref{rev:ex-affine-OU} with a unique inverse and exact Jacobian.\\
6 & Exponential\allowbreak{}TransformedOU & C & Example~\ref{rev:ex-exponential-OU}, conditional on support and the intended target $m$; the original zero target must not be interpreted automatically as one dollar.\\
7 & SpecialFunction\allowbreak{}TransformedOU & N & Proposition~\ref{rev:prop-OU-Kummer-failure}; a forward map and derivative do not constitute an inverse transition density.\\
8 & AffineTransformedCIR & S & Example~\ref{rev:ex-affine-CIR}, subject to CIR boundary and observation-support checks.\\
9 & Quadratic\allowbreak{}TransformedCIR & C & Proposition~\ref{rev:prop-quadratic-CIR}; only a stated finite window satisfying \eqref{rev:quadratic-monotonicity}.\\
10 & Exponential\allowbreak{}TransformedCIR & C & Example~\ref{rev:ex-exponential-CIR}; enforce \eqref{rev:exp-CIR-moment} under stationary start and the exact lower support \eqref{rev:exp-CIR-support}.\\
11 & SpecialFunction\allowbreak{}TransformedCIR & C & Proposition~\ref{rev:prop-CIR-Kummer}; a specified branch must pass monotonicity, inverse, support, moment, and numerical-stability checks.  A finite branch scan is diagnostic, not a full density implementation.\\
12 & MOUF & S & Constant-intensity stationary OU reset; Theorems~\ref{rev:thm-mixed-kernel}--\ref{rev:thm-poisson-construction}.\\
13 & MCIR & S & Stationary CIR reset under the full-density and no-boundary-atom conditions of Proposition~\ref{rev:prop-diffusion-copula-generator}; a base-density fallback and weak identification are empirical qualifications, not changes to the density-mixture theorem.\\
14 & SeasonalIntensity\allowbreak{}MOU & S & The nonhomogeneous Poisson version of Theorem~\ref{rev:thm-poisson-construction}, provided the nonnegative seasonal intensity is right-continuous and locally bounded.\\
15 & JumpOU & B & An exact compound-Poisson OU with $\sigma>0$ is an absolutely continuous jump-diffusion extension; the pure-jump boundary $\sigma=0$ has a no-jump atom and lies outside the present Copula-density scope.  With jump mean $\mu_J$, its mean target is $\bar r+\lambda\mu_J/\kappa$.  The current one-jump moment approximation is not a proved Chapman--Kolmogorov semigroup.\\
16 & MarkovSwitching\allowbreak{}OU & L & $(X_t,S_t)$ is Markov; $X_t$ alone generally is not.  Structural claims must be stated on the augmented regime state.\\
17 & NestedOU & L & The vector of two OU factors is Markov.  Its scalar sum is Markov only in degenerate/equal-speed cases; pairwise Gaussian composite scores are not a full Kalman path likelihood.\\
18 & RandomWalk & B & Explicit nonreverting negative control.\\
19 & GaussianGamma\allowbreak{}Stationary & N & It is a valid stationary Markov density, but it is not a proved mean-reverting model until \eqref{rev:gamma-full-sign} passes; the scalar root is insufficient.\\
20 & GaussianGamma\allowbreak{}TimeVarying & N & The endpoint ODE \eqref{rev:gamma-time-necessary-ode} is only diagnostic.  A scale path, full density, and global sign test \eqref{rev:gamma-time-full-sign} are still required.\\
21 & ParametricCopula\allowbreak{}Grid & L & Gaussian $\rho(h)=e^{-\kappa h}$ and its independence-reset mixture have continuous-time semigroups.  At admissible interior parameters, Student-$t$, Clayton, Gumbel, and Frank cells give absolutely continuous one-step Copula densities but lack a proved same-family continuous-time embedding.  Singular boundary parameters are excluded.  Iterating a fixed five-minute Copula-density transition operator defines a discrete-time Markov chain, not the claimed parametric semigroup.\\
22 & TwoScale\allowbreak{}Gaussian & L & A two-OU latent vector is valid; the scalar correlation mixture is generally non-Markov and lies outside the one-dimensional generator theorem.\\
23 & TwoScale\allowbreak{}Mixed & L & A rigorous model resets the entire latent vector at Poisson times.  Pairwise scalar mixed densities alone do not prove the density Chapman--Kolmogorov identity or process-level consistency.\\
24 & FullyTimeVarying\allowbreak{}MarginalCopula & C & If $X_t=\mu(t)+s(t)Z_t$ with standardized OU $Z$ and deterministic $s(t)>0$, then $X$ itself is a one-dimensional time-inhomogeneous Markov process with exact Gaussian transition and Copula densities.  Reversion to fixed $m$ additionally requires $\beta(t)=\kappa-s'(t)/s(t)>0$, $\mu'(t)=\beta(t)\{m-\mu(t)\}$, and $\mu(t)\to m$.  Permanent seasonality or a nonzero linear trend violates the limiting-mean condition, not Markovity.\\
25 & MixedCopula\allowbreak{}BondPDE & S & Theorem~\ref{rev:thm-bond-PDE} under $\Qq$.  It is not identifiable from a physical-measure USDT spot series without bond-price data.\\
\bottomrule
\end{longtable}
\normalsize

\subsection{Further formulas for the extension classes}

For a two-sided threshold diffusion centered at $c$, with drift
$-\kappa_+(x-c)$ above $c$ and $-\kappa_-(x-c)$ below $c$, and piecewise
constant volatilities $\sigma_\pm$, its invariant density is proportional to
\begin{equation}
\pi(x)\propto
\begin{cases}
\sigma_+^{-2}\exp\{-\kappa_+(x-c)^2/\sigma_+^2\},&x\ge c,\\
\sigma_-^{-2}\exp\{-\kappa_-(x-c)^2/\sigma_-^2\},&x<c.
\end{cases}
\label{rev:threshold-invariant}
\end{equation}
Direct integration shows $\int(x-c)\pi(x)\dd x=0$ if and only if
$\kappa_+=\kappa_-$.  Thus directional drift toward $c$ is not enough to call
$c$ the asymptotic mean.

For a time-varying Gaussian-margin model
$X_t=\mu(t)+s(t)Z_t$, where
$\dd Z_t=-\kappa Z_t\dd t+\sqrt{2\kappa}\dd W_t$, It\^o's formula gives
\begin{equation}
b_X(t,x)=\mu'(t)+\left\{\frac{s'(t)}{s(t)}-\kappa\right\}
\{x-\mu(t)\}.
\label{rev:timevarying-normal-drift}
\end{equation}
Writing $\beta(t)=\kappa-s'(t)/s(t)$, this equals
$-\beta(t)(x-m)$ precisely when
\begin{equation}
\beta(t)>0,\qquad
\mu'(t)=\beta(t)\{m-\mu(t)\}.
\label{rev:timevarying-normal-MR-condition}
\end{equation}

For an OU jump diffusion with compound-Poisson jump mean $\mu_J$ and intensity
$\lambda$, the forward drift is
\begin{equation}
\kappa(\bar r-x)+\lambda\mu_J
=\kappa\left(\bar r+\frac{\lambda\mu_J}{\kappa}-x\right).
\label{rev:jump-OU-target}
\end{equation}
Hence using the base OU level as the target is valid only for centered jumps.

\section{What the corrected theory does and does not claim}

\begin{enumerate}[label=(\roman*)]
\item Exact finite-time identities such as \eqref{rev:exact-mr-copula} take
priority over a formal local PDE.  A zero generator drift is not silently
upgraded from local to true martingale.
\item A formal PDE solution becomes an observable one-dimensional Markov model
only after strict monotonicity, support, inverse-Jacobian, moment, and boundary
conditions pass.
\item The stationary and time-varying Gamma equations are retained as necessary
diagnostics, not existence theorems.
\item The quadratic CIR transform is a finite-horizon model.  The OU even
special-function transform is excluded from the monotone framework.  The CIR
Kummer family is branch-specific, and the erroneous second CIR reduction is
replaced by \eqref{rev:CIR-correct-reduction}.
\item The absolutely continuous mixed-copula reset construction is rigorous and gives the exact
structural improvement \eqref{rev:mixed-exact-mr}; empirical superiority still
has to be demonstrated by common out-of-sample density and calibration metrics.
\item Hidden-regime, two-factor, and two-scale specifications belong to an
augmented-state theory.  Horizon-specific copula fits and transition
approximations are prediction models unless a consistent continuous-time
semigroup is separately proved.
\end{enumerate}

\fi
\clearpage
\appendix
\section{Detailed proofs}
\label{rev:sec-detailed-proofs}

This appendix expands the measure-theoretic steps that are easy to hide inside
phrases such as ``interchange the limit and the conditional expectation.''
It covers the general results used repeatedly in the main text.  The
model-specific OU, CIR, and special-function calculations remain next to their
statements because they are direct substitutions into the displayed
differential equations.

\subsection{Copula kernels, density versions, and partial derivatives}
\label{rev:app-copula-kernel-proof}

\begin{proof}[Detailed proof of the kernel and derivative claims in
Proposition~\ref{rev:prop-density-copula-consistency}]
Fix $s<t$ and suppress the time subscripts temporarily.  For the selected
jointly Borel density $c$, define
\[
K(u,B)=\int_Bc(u,v)\dd v,
\qquad
H^c(u,y)=K(u,(0,y])=\int_0^yc(u,v)\dd v.
\]
Row normalization makes $B\mapsto K(u,B)$ a probability measure for every
$u$.  Joint measurability of $c$, first for intervals and then by the monotone
class theorem for all $B\in\mathcal B(I)$, makes $u\mapsto K(u,B)$ Borel.
Column normalization and Tonelli's theorem give, for every Borel $B$,
\begin{align}
\int_IK(u,B)\dd u
&=\int_B\left\{\int_Ic(u,v)\dd u\right\}\dd v
=\Leb(B),
\label{rev:app-kernel-invariance}
\end{align}
so the selected kernel preserves the uniform law.

For every $x,y\in[0,1]$, the copula reconstructed from $c$ satisfies
\begin{equation}
C(x,y)=\int_0^xH^c(a,y)\dd a.
\label{rev:app-copula-kernel-reconstruction}
\end{equation}
For fixed $y$, the function $H^c(\cdot,y)$ is measurable and bounded by one.
The fundamental theorem for Lebesgue integrals therefore implies
\begin{equation}
\partial_1C(u,y)=H^c(u,y)
\quad\text{for $\Leb$-a.e. }u.
\label{rev:app-partial-from-kernel}
\end{equation}
This is an equality for each fixed $y$; its exceptional set may initially
depend on $y$.  Because $H^c(u,\cdot)$ and every conditional-CDF version are
right-continuous, equality at all rational $y$ outside one null set extends to
all $y\in[0,1]$ outside that same source-state null set.

At any source state $u$ where $a\mapsto H^c(a,y)$ is continuous, the pointwise
upgrade follows directly from the integral difference quotient:
\begin{align}
\frac{C(u+h,y)-C(u,y)}{h}-H^c(u,y)
&=\frac1h\int_u^{u+h}\{H^c(a,y)-H^c(u,y)\}\dd a,\notag\\
\left|\frac{C(u+h,y)-C(u,y)}{h}-H^c(u,y)\right|
&\le \sup_{|a-u|\le |h|}|H^c(a,y)-H^c(u,y)|
\longrightarrow0.
\label{rev:app-pointwise-partial-upgrade}
\end{align}
Thus a continuity condition on the integrated density row, not merely existence
of a $\Leb^2$-a.e. density, is what upgrades
\eqref{rev:app-partial-from-kernel} to a pointwise derivative identity.  The
modified partial Dini construction of \citet{FangJiangLiuYang2020} supplies
another pointwise CDF-valued version $\mathfrak D_1C(u,\cdot)$ satisfying the
same reconstruction formula.  Applying
\eqref{rev:app-partial-from-kernel} to both versions shows that
$\mathfrak D_1C(u,\cdot)=H^c(u,\cdot)$ for a.e. $u$, although equality can fail
on a source-null set.  These uses of absolute continuity and differentiation
are instances of the results in \citet[Chapter~3]{Folland1999}.

Now fix $0\le s<r<t$ and $u\in I$.  Since
\eqref{rev:copula-density-CK} holds for a.e. $v$, equality of its integrals over
an arbitrary $B\in\mathcal B(I)$ and Tonelli's theorem give
\begin{align}
K_{s,t}(u,B)
&=\int_B\int_Ic_{s,r}(u,w)c_{r,t}(w,v)\dd w\dd v\notag\\
&=\int_Ic_{s,r}(u,w)
       \left\{\int_Bc_{r,t}(w,v)\dd v\right\}\dd w\notag\\
&=\int_IK_{r,t}(w,B)K_{s,r}(u,\dd w).
\label{rev:app-density-to-kernel-CK}
\end{align}
This proves the kernel Chapman--Kolmogorov equation for every selected source
$u$ and every Borel $B$.  Conversely, if both sides are absolutely continuous
in the target variable, equality for every $B$ identifies their
Radon--Nikodym densities only for a.e. $v$.

Finally, let $C=C_{s,r}$ and $D=C_{r,t}$.  The functions
\[
w\longmapsto\int_0^xc_{s,r}(a,w)\dd a,
\qquad
w\longmapsto\int_0^yc_{r,t}(w,b)\dd b
\]
are measurable a.e. representatives of $\partial_2C(x,w)$ and
$\partial_1D(w,y)$, respectively.  Substitution into
\eqref{rev:star-product-density-version}, followed by Tonelli and density CK,
gives
\begin{align}
(C_{s,r}*C_{r,t})(x,y)
&=\int_I\int_0^xc_{s,r}(a,w)\dd a
             \int_0^yc_{r,t}(w,b)\dd b\dd w\notag\\
&=\int_0^x\int_0^y
     \left\{\int_Ic_{s,r}(a,w)c_{r,t}(w,b)\dd w\right\}\dd b\dd a\notag\\
&=\int_0^x\int_0^yc_{s,t}(a,b)\dd b\dd a
=C_{s,t}(x,y).
\label{rev:app-star-product-integral}
\end{align}
All changes in integration order are justified by nonnegativity.  Conversely,
equality of these copula CDFs identifies their two-dimensional densities only
$\Leb^2$-a.e. and hence, by Fubini, yields conditional density composition only
for a.e. source state.  This proves why star-product consistency alone is weaker
than the selected pointwise kernel condition in
Assumption~\ref{rev:ass-copula-density}.
\end{proof}

\subsection{Short-time kernels acting on moving functions}

\begin{lemma}[Moving-function kernel lemma]
\label{rev:app-kernel-lemma}
Fix $t\ge0$ and an interior state $u\in I$.  Write
\[
K_h(u,\dd v)=c_{t,t+h}(u,v)\dd v,
\qquad h>0,
\]
so that $K_h(u,I)=1$.  The following statements hold.
\begin{enumerate}[label=(\alph*)]
\item Suppose that
\begin{equation}
K_h\bigl(u,\{v:|v-u|\ge\delta\}\bigr)\longrightarrow0
\quad\text{for every }\delta>0.
\label{rev:kernel-stochastic-continuity}
\end{equation}
If $g$ is bounded and continuous at $u$, then
\begin{equation}
\int_Ig(v)K_h(u,\dd v)\longrightarrow g(u).
\label{rev:kernel-fixed-function-limit}
\end{equation}
\item In addition to part~(a), let $r_h:I\to\R$ be measurable and suppose
$g$ is bounded and continuous at $u$ and
\begin{equation}
\lVert r_h-g\rVert_\infty\longrightarrow0.
\label{rev:moving-uniform-limit}
\end{equation}
Then
\begin{equation}
\int_Ir_h(v)K_h(u,\dd v)\longrightarrow g(u).
\label{rev:moving-kernel-limit}
\end{equation}
\item For an unbounded $g$, let $w:I\to[1,\infty)$ be measurable and define
$\lVert f\rVert_w=\sup_{v\in I}|f(v)|/w(v)$.  If
\begin{align}
\lVert r_h-g\rVert_w&\longrightarrow0,
\label{rev:moving-weighted-limit}\\
\sup_{0<h<h_0}\int_Iw(v)K_h(u,\dd v)&<\infty,
\label{rev:kernel-weight-bound}\\
\int_Ig(v)K_h(u,\dd v)&\longrightarrow g(u),
\label{rev:weighted-strong-continuity}
\end{align}
then \eqref{rev:moving-kernel-limit} still holds.
\end{enumerate}
\end{lemma}

\begin{proof}
For part~(a), split the integral at distance $\delta$ from $u$:
\begin{align*}
\left|\int_Ig(v)K_h(u,\dd v)-g(u)\right|
&\le \int_I|g(v)-g(u)|K_h(u,\dd v)\\
&\le \sup_{|v-u|<\delta}|g(v)-g(u)|
 +2\lVert g\rVert_\infty
 K_h\bigl(u,\{|v-u|\ge\delta\}\bigr).
\end{align*}
First let $h\downarrow0$ and use
\eqref{rev:kernel-stochastic-continuity}; then let $\delta\downarrow0$ and
use continuity of $g$ at $u$.  This proves
\eqref{rev:kernel-fixed-function-limit}, or equivalently
$K_h(u,\cdot)\Rightarrow\delta_u$.

For part~(b), row normalization gives
\begin{align*}
\left|\int_Ir_h(v)K_h(u,\dd v)-g(u)\right|
&\le \int_I|r_h(v)-g(v)|K_h(u,\dd v)\\
&\quad+\left|\int_Ig(v)K_h(u,\dd v)-g(u)\right|\\
&\le \lVert r_h-g\rVert_\infty
 +\left|\int_Ig(v)K_h(u,\dd v)-g(u)\right|.
\end{align*}
Both terms tend to zero by \eqref{rev:moving-uniform-limit} and part~(a).

For part~(c), replace the first term in the preceding bound by
\[
\int_I|r_h(v)-g(v)|K_h(u,\dd v)
\le \lVert r_h-g\rVert_w\int_Iw(v)K_h(u,\dd v).
\]
Equations \eqref{rev:moving-weighted-limit}--
\eqref{rev:weighted-strong-continuity} prove the claim.  This weighted version
is the relevant one for Normal, Student, or other quantiles that diverge at the
endpoints of $I$.
\end{proof}

Pointwise convergence $r_h(v)\to g(v)$ alone is not enough in this lemma: the
measure $K_h(u,\dd v)$ also changes with $h$.  Conditions
\eqref{rev:moving-uniform-limit} or
\eqref{rev:moving-weighted-limit}--\eqref{rev:kernel-weight-bound} provide the
uniform integrability needed for the moving integrand.

\subsection{Full integral proof of the quantile--generator drift identity}
\label{rev:app-drift-identity-proof}

\begin{proof}[Detailed proof of Proposition~\ref{rev:prop-generator-sign}]
Fix $t$ and put $q_t(v)=q(t,v)=Q_t(v)$ and
$\dot q_t(v)=\partial_tq(t,v)$.  Assume that all displayed first moments are
finite and that, for the state $u$ under consideration,
\begin{equation}
\mathcal A_tq_t(u)
=\lim_{h\downarrow0}
\frac{\int_Iq_t(v)c_{t,t+h}(u,v)\dd v-q_t(u)}{h}
\label{rev:app-generator-integral}
\end{equation}
exists.  Also assume either part~(b) or part~(c) of
Lemma~\ref{rev:app-kernel-lemma} with
\begin{equation}
r_h(v)=\frac{q(t+h,v)-q(t,v)}{h},
\qquad g(v)=\dot q_t(v).
\label{rev:app-time-difference-quotient}
\end{equation}
These are explicit sufficient conditions for the limit-interchange clause in
Proposition~\ref{rev:prop-generator-sign}.

For a regular conditional kernel version, strict invertibility of $q_t$ gives,
for $\Leb$-almost every $u$,
\begin{equation}
P^X_{t,t+h}\bigl(q_t(u),B\bigr)
=\int_I\one_{B}\bigl(q_{t+h}(v)\bigr)c_{t,t+h}(u,v)\dd v.
\label{rev:app-pushforward-kernel}
\end{equation}
Consequently,
\begin{align}
b_X(t,q_t(u))
&=\lim_{h\downarrow0}
\frac{\int_Iq_{t+h}(v)c_{t,t+h}(u,v)\dd v-q_t(u)}{h}.
\label{rev:app-forward-drift-integral}
\end{align}
Using $\int_Ic_{t,t+h}(u,v)\dd v=1$, the difference quotient in
\eqref{rev:app-forward-drift-integral} is exactly
\begin{align}
&\int_I
\frac{q_{t+h}(v)-q_t(v)}{h}c_{t,t+h}(u,v)\dd v
\label{rev:app-drift-time-term}\\
&\qquad+
\frac{\int_Iq_t(v)c_{t,t+h}(u,v)\dd v-q_t(u)}{h}.
\label{rev:app-drift-state-term}
\end{align}
Lemma~\ref{rev:app-kernel-lemma} and
\eqref{rev:app-time-difference-quotient} show that
\eqref{rev:app-drift-time-term} tends to $\dot q_t(u)$.  Definition
\eqref{rev:app-generator-integral} shows that
\eqref{rev:app-drift-state-term} tends to $\mathcal A_tq_t(u)$.  Hence
\begin{equation}
b_X(t,Q_t(u))
=\partial_tq(t,u)+\mathcal A_tq(t,\cdot)(u)
=(\partial_t+\mathcal A_t)q(t,u).
\label{rev:app-full-drift-identity}
\end{equation}
Because $U_t$ is uniform, the statement is naturally a
$\Leb$-almost-everywhere statement in $u$.  A continuous kernel version extends
it to every interior state where the limits exist.  Strict monotonicity of
$q_t$ then transfers the sign of the right-hand side to the physical state
$x=q_t(u)$, which proves the sign assertion in
Proposition~\ref{rev:prop-generator-sign} after the limiting-mean condition is
imposed separately.
\end{proof}

\subsection{Conditional expectations and observable transition densities}
\label{rev:app-conditional-kernel-proof}

\begin{proof}[Detailed proof of Theorems~\ref{rev:thm-exact-density} and
\ref{rev:thm-exact-mr}, and Proposition~\ref{rev:prop-observable-density}]
For $s<t$, define
\[
H_{s,t}(u)=\int_Iq_t(v)c_{s,t}(u,v)\dd v.
\]
Column normalization and Tonelli's theorem give
\begin{align*}
\int_I|H_{s,t}(u)|\dd u
&\le\int_I\int_I|q_t(v)|c_{s,t}(u,v)\dd v\dd u\\
&=\int_I|q_t(v)|
   \left\{\int_Ic_{s,t}(u,v)\dd u\right\}\dd v\\
&=\int_I|q_t(v)|\dd v<\infty.
\end{align*}
For a bounded $\mathcal F_s$-measurable random variable $Z$, the Markov
property, first for bounded truncations of $q_t$ and then by $L^1$ convergence,
gives
\begin{align*}
\E[ZX_t]
&=\E[Zq_t(U_t)]\\
&=\E\!\left[Z\int_Iq_t(v)c_{s,t}(U_s,v)\dd v\right]
=\E[ZH_{s,t}(U_s)].
\end{align*}
This is the defining test-function identity for
\[
\E[X_t\mid\mathcal F_s]=H_{s,t}(U_s).
\]
Since $X_s=q_s(U_s)$ and $U_s\sim\Leb$, equality
$H_{s,t}(U_s)=q_s(U_s)$ almost surely is equivalent to
$H_{s,t}(u)=q_s(u)$ for $\Leb$-almost every $u$.  This proves both directions
of Theorem~\ref{rev:thm-exact-density} for the fixed pair $(s,t)$; no common
null set over the uncountable collection of time pairs is asserted.  Replacing
the last equality by
\[
H_{s,t}(u)-m=\mathcal R_{s,t}\{q_s(u)-m\}
\]
proves both directions of Theorem~\ref{rev:thm-exact-mr}.

Under the density assumptions of
Proposition~\ref{rev:prop-observable-density}, the substitution
$v=F_t(y)$, $\dd v=p_t(y)\dd y$, in the conditional kernel gives
\[
p^X_{s,t}(x,y)
=c_{s,t}\{F_s(x),F_t(y)\}p_t(y).
\]
The three required identities follow by explicit changes of variables:
\begin{align*}
\int_{E_t^\circ}p^X_{s,t}(x,y)\dd y
&=\int_Ic_{s,t}\{F_s(x),v\}\dd v=1,\\
\int_{E_s^\circ}p_s(x)p^X_{s,t}(x,y)\dd x
&=p_t(y)\int_Ic_{s,t}\{u,F_t(y)\}\dd u=p_t(y),\\
\int_{E_r^\circ}p^X_{s,r}(x,z)p^X_{r,t}(z,y)\dd z
&=p_t(y)\int_Ic_{s,r}\{F_s(x),w\}
                   c_{r,t}\{w,F_t(y)\}\dd w\\
&=p_t(y)c_{s,t}\{F_s(x),F_t(y)\}
=p^X_{s,t}(x,y).
\end{align*}
The first line holds for every selected $x$, the second for Lebesgue-a.e. $y$,
and the third for every selected $x$ and Lebesgue-a.e. $y$.  The final equality
is precisely the density Chapman--Kolmogorov identity; integrating it over an
arbitrary target Borel set yields the corresponding pointwise kernel identity.
\end{proof}

\subsection{The generator martingale theorem and rescaling PDE}
\label{rev:app-generator-martingale-proof}

\begin{proof}[Detailed proof of Theorem~\ref{rev:thm-generator-martingale}]
Write
\[
B(t,u)=(\partial_t+\mathcal A_t)q(t,u).
\]
Membership in the time--space extended generator domain means that, after a
localizing sequence if necessary,
\begin{equation}
L_t=q(t,U_t)-q(0,U_0)-\int_0^tB(s,U_s)\dd s
\label{rev:app-dynkin-local-martingale}
\end{equation}
is a local martingale and the finite-variation integral is locally finite.
If $B=0$ for $\dd t\otimes\Leb(\dd u)$-almost every $(t,u)$, uniformity of
$U_s$ and Tonelli's theorem yield, for every $T<\infty$,
\begin{align*}
\E\int_0^T\one_{\{B(s,U_s)\ne0\}}\dd s
&=\int_0^T\int_I\one_{\{B(s,u)\ne0\}}\dd u\dd s=0.
\end{align*}
Thus the integral in \eqref{rev:app-dynkin-local-martingale} vanishes and
$q(t,U_t)$ is a local martingale.  Class--$(D)$ on each finite horizon makes
the stopped family uniformly integrable, so the local martingale is a true
martingale.

Conversely, suppose $q(t,U_t)$ is a true martingale.  It is also a local
martingale, and subtracting the local martingale $L$ in
\eqref{rev:app-dynkin-local-martingale} shows that
\[
A_t:=\int_0^tB(s,U_s)\dd s
\]
is a local martingale.  It is an absolutely continuous finite-variation
process starting from zero.  A finite-variation local martingale is
indistinguishable from a constant, so $A_t=0$ for all $t$ almost surely and
$B(s,U_s)=0$ for $\dd s\otimes\Pp$-almost every $(s,\omega)$.  Uniformity of
$U_s$ and Tonelli's theorem now give
\[
0=\int_0^T\Pp\{B(s,U_s)\ne0\}\dd s
=\int_0^T\int_I\one_{\{B(s,u)\ne0\}}\dd u\dd s.
\]
Hence $B=0$ for $\dd t\otimes\Leb(\dd u)$-almost every $(t,u)$, proving the
converse in Theorem~\ref{rev:thm-generator-martingale}.
\end{proof}

\begin{proof}[Detailed generator proof of
Proposition~\ref{rev:prop-martingale-rescaling}]
For completeness, define the deterministic scaling factor
\[
R_t=\mathcal R_{0,t}=\exp\!\left\{\int_0^t\gamma(a)\dd a\right\},
\]
and set
\[
\widetilde q(t,u)=\frac{q(t,u)-m}{R_t}.
\]
Since $R_t$ is deterministic, $R_t'=\gamma(t)R_t$, and
$\mathcal A_t$ acts only on the state variable,
\begin{align}
(\partial_t+\mathcal A_t)\widetilde q(t,u)
=\frac{1}{R_t}
\left[(\partial_t+\mathcal A_t)q(t,u)
-\gamma(t)\{q(t,u)-m\}\right].
\label{rev:app-rescaled-generator}
\end{align}
Thus the rescaled process is a martingale, subject to the preceding domain and
class--$(D)$ qualifications, exactly when
\eqref{rev:strong-mr-generator-PDE} holds.  At the finite-horizon level,
multiplicativity gives $R_t=R_s\mathcal R_{s,t}$ and therefore
\begin{align*}
\E\!\left[\frac{X_t-m}{R_t}\middle|\mathcal F_s\right]
&=\frac{\mathcal R_{s,t}(X_s-m)}{R_s\mathcal R_{s,t}}
=\frac{X_s-m}{R_s}.
\end{align*}
The reverse calculation proves the converse part of
Proposition~\ref{rev:prop-martingale-rescaling} without appealing only to formal PDE
algebra.
\end{proof}

\subsection{Diffusion transforms and reset mixtures}
\label{rev:app-transform-reset-proofs}

\begin{proof}[Detailed proof of Theorem~\ref{rev:thm-diffusion-transform}]
For Theorem~\ref{rev:thm-diffusion-transform}, It\^o's formula gives
\begin{align*}
\dd\{f(t,R_t)-m\}
&=(\partial_t+\mathcal L)f(t,R_t)\dd t
  +\sigma(R_t)f_x(t,R_t)\dd W_t\\
&=\gamma(t)\{f(t,R_t)-m\}\dd t
  +\sigma(R_t)f_x(t,R_t)\dd W_t.
\end{align*}
With $R_t^\gamma=\exp\{\int_0^t\gamma(a)\dd a\}$, the deterministic product
rule yields
\begin{equation}
\dd\!\left[\frac{f(t,R_t)-m}{R_t^\gamma}\right]
=\frac{\sigma(R_t)f_x(t,R_t)}{R_t^\gamma}\dd W_t.
\label{rev:app-transform-martingale}
\end{equation}
Condition \eqref{rev:transform-integrability} makes the stochastic integral in
\eqref{rev:app-transform-martingale} square integrable on every finite
horizon.  Taking the conditional expectation between $s$ and $t$ and using
$R_t^\gamma/R_s^\gamma=\mathcal R_{s,t}$ gives
\eqref{rev:transform-conditional-mean}.  Strict monotonicity gives the
time-dependent inverse $R_s=f^{-1}(s,Y_s)$ and the transition kernel
\[
P^Y_{s,t}(y,B)
=P^R_{s,t}\bigl(f^{-1}(s,y),f^{-1}(t,B)\bigr),
\]
which proves the Markov claim.  Conversely, the conditional-mean identity
makes the left side of \eqref{rev:app-transform-martingale} a martingale;
uniqueness of the finite-variation term in It\^o's decomposition forces
\eqref{rev:transform-PDE} wherever the coefficients and derivatives are
defined.
\end{proof}

\begin{proof}[Detailed proof of Theorems~\ref{rev:thm-mixed-kernel} and
\ref{rev:thm-mixed-mr}]
Finally, write $a=a_{s,r}$ and $b=a_{r,t}$.  Invariance gives
$P_{s,r}\Pi=\Pi P_{r,t}=\Pi$ and $\Pi^2=\Pi$, so the reset-kernel composition
in Theorem~\ref{rev:thm-mixed-kernel} expands as
\begin{align*}
&\{aP_{s,r}+(1-a)\Pi\}
 \{bP_{r,t}+(1-b)\Pi\}\\
&\quad=abP_{s,t}
 +\{a(1-b)+(1-a)b+(1-a)(1-b)\}\Pi\\
&\quad=abP_{s,t}+(1-ab)\Pi
=P^{\mathrm{mix}}_{s,t},
\end{align*}
because $ab=a_{s,t}$.  For the Copula densities the corresponding integral is
\begin{align*}
&\int_I\{ac_{s,r}(u,w)+1-a\}
          \{bc_{r,t}(w,v)+1-b\}\dd w\\
&\quad=ab c_{s,t}(u,v)+a(1-b)+(1-a)b+(1-a)(1-b)\\
&\quad=a_{s,t}c_{s,t}(u,v)+1-a_{s,t},
\end{align*}
where row normalization, column normalization, and the base
Chapman--Kolmogorov identity were used term by term.  This density equality is
for every selected $u$ and a.e. $v$; its integrated kernel equality is for every
$u$ and every Borel target set.  Moreover,
$a_{t,t+h}=1-\lambda(t)h+o(h)$ and
$P_{t,t+h}h_0=h_0+h\mathcal G_th_0+o(h)$ on the generator domain.  Therefore
\begin{align*}
\frac{P^{\mathrm{mix}}_{t,t+h}h_0-h_0}{h}
&=\mathcal G_th_0
 +\lambda(t)\{\Pi h_0-h_0\}+o(1),
\end{align*}
which proves \eqref{rev:mixed-generator}.  The conditional centered first
moment is, directly in integral form,
\begin{align*}
\int_E(y-m)P^{\mathrm{mix}}_{s,t}(x,\dd y)
&=a_{s,t}\int_E(y-m)P_{s,t}(x,\dd y)\\
&\quad+(1-a_{s,t})\int_E(y-m)\pi(\dd y)\\
&=a_{s,t}\mathcal R_{s,t}(x-m),
\end{align*}
and its right derivative at $t=s$ gives
$b_{\mathrm{mix}}(s,x)=b_0(s,x)+\lambda(s)(m-x)$.
\end{proof}

\endinput
\clearpage
\bibliographystyle{plainnat}
\bibliography{meanreversion_theory_references}

\end{document}
